% ============================================================================
%  Dirichlet-Charakter-Atlas: Eine Weil-Kern-Kartografie
%  von Low-Conductor-Sektor-Gaps
%
%  Autor: Lukas Geiger, Independent Researcher, Bernau
%  Entwurf v2 — überarbeitet nach N=600-Serveranalyse
%  Status: Finaler Entwurf
% ============================================================================
\documentclass[a4paper,10pt,twocolumn]{article}

\usepackage[a4paper,top=2.0cm,bottom=2.5cm,left=1.8cm,right=1.8cm,columnsep=0.6cm]{geometry}
\usepackage[utf8]{inputenc}
\usepackage[T1]{fontenc}
\usepackage{lmodern}
\usepackage{microtype}
\usepackage[ngerman]{babel}
\usepackage{amsmath,amssymb,amsfonts,mathtools,bm}
\usepackage{graphicx,float,booktabs,array,multirow}
\usepackage{xcolor}
\usepackage{xurl}
\usepackage[unicode=true,colorlinks=true,linkcolor=blue!60!black,citecolor=blue!60!black,urlcolor=blue!60!black]{hyperref}
\usepackage[numbers,square,sort&compress]{natbib}
\bibliographystyle{unsrtnat}

\usepackage{titlesec}
\renewcommand{\thesection}{\Roman{section}}
\renewcommand{\thesubsection}{\Alph{subsection}}

\usepackage{amsthm}
\newtheorem{theorem}{Theorem}[section]
\newtheorem{proposition}[theorem]{Proposition}
\newtheorem{corollary}[theorem]{Corollary}
\newtheorem{lemma}[theorem]{Lemma}
\theoremstyle{definition}
\newtheorem{definition}[theorem]{Definition}
\newtheorem{example}[theorem]{Beispiel}
\theoremstyle{remark}
\newtheorem{remark}[theorem]{Bemerkung}

% --- Eigene Makros ---
\newcommand{\gap}{\operatorname{gap}}
\newcommand{\chifn}{\chi}
\newcommand{\Zmodq}{(\mathbb{Z}/q\mathbb{Z})^*}
\newcommand{\halfline}{\tfrac{1}{2}}
\newcommand{\Deltachi}{\Delta_{\chifn}}
\newcommand{\archdiff}{\operatorname{arch\_diff}}
\newcommand{\primediff}{\operatorname{prime\_diff}_{\mathrm{matrix}}}
\DeclareMathOperator{\sgn}{sgn}
\providecommand*{\HyperDestNameFilter}[1]{de.#1}

% ============================================================================
\title{Der Dirichlet-Charakter-Atlas\\
  \large --- Ein Weil-Kern-Zahlenatlas für Sektor-Gaps von Dirichlet-$L$-Funktionen\\
  \normalsize Galerkin-Diagnostik und die Grenzen der Theorie führender Ordnung}
\author{Lukas Geiger\\\small Independent Researcher, Bernau}
\date{\today}

\hypersetup{
  pdftitle={Der Dirichlet-Charakter-Atlas},
  pdfauthor={Lukas Geiger},
  pdfsubject={Weil-Kern-Zahlenatlas für Sektor-Gaps von Dirichlet-L-Funktionen},
  pdfkeywords={Dirichlet-L-Funktionen, Weilsche quadratische Form, Sektor-Gaps, Galerkin-Diagnostik, Hilbert-P\'olya}
}

\begin{document}
\twocolumn[
\maketitle
\begin{abstract}
\noindent
Diese Arbeit ist eine Instanziierung der Funktionalen Stabilitätstheorie
(Functional Stability Theory, FST) im mathematischen Bereich: Sie erweitert die
Dreikomponenten-Zerlegung $\mathrm{RH}(X) = \mathrm{GBC}(X) + C^{(2)}(X) + \mathrm{Hurwitz}(X)$
(siehe Math-Master~\cite{MathMaster}, \S2.4) auf
getwistete $L$-Funktionen, wobei der Sektor-Gap (sector gap) die quantitative
Signatur von $C^{(2)}_{\chifn}$ darstellt. Der Math-Master
\cite{MathMaster} entwickelt die übergreifende Zeta-Zoo-Architektur und
die Trinität aus UCU + SGE + Weil-Form-Transversalität; der Physics-Master
\cite{RFEPFramework} entwickelt die physikalische Begleitformulierung
über das Renormierte Prinzip der Freien Energie (Renormalized Free Energy Principle).

Wir entwickeln einen numerischen Weil-Kern-Atlas von Sektor-Gaps für Dirichlet
$L$-Funktionen, die zu zehn reellen Charakteren mit kleinem Führer
$\chi_D$, $D\in\{5,8,12,13,17,21,24,29,33,60\}$, assoziiert sind, bei Galerkin-Trunkierungen
$N\in\{200,400,600\}$ und einem Cut-off von $\lambda=20000$. Innerhalb der
Dreikomponenten-Zerlegung ist der Sektor-Gap die quantitative
Signatur von $C^{(2)}_{\chifn}$.
Unsere Hauptergebnisse sind \emph{negativ, aber strukturell informativ}. Erstens
wird die Approximation führender Ordnung $\phi^+_{\chifn}\approx\phi^-_{\chifn}$
der zwei Sektor-Grundzustände empirisch falsifiziert: Ein systematischer
16-Fenster-Test, der auch den empirischen Grundzustand selbst einschließt, verfehlt das
Erfolgskriterium von $9/10$ Vorzeichengenauigkeit (sign accuracy), $R^2\geq 0.8$. Zweitens
reduziert sich die scheinbar „exakte“ Formel
\[
  \gap_{\chifn} = -2\!\!\sum_{p,m}\!\!\frac{\chifn(p)^m\log p}{p^{m/2}}
    \Deltachi(m\log p) + \archdiff_{\chifn} + O(N^{-1})
\]
mit
$\Deltachi(t):=(\phi^+_{\chifn}\ast\phi^+_{\chifn})(t)+(\phi^-_{\chifn}\ast\phi^-_{\chifn})(t)$,
sobald beide Grundzustände aus derselben Galerkin-Diagonalisierung entnommen werden,
zu einer Matrixzerlegungs-Tautologie: Eine scheinbare
$9/9$ Vorzeichengenauigkeit bei $N=200$ verschwindet als numerisches Interpolations-Artefakt,
wenn die Analyse bei $N=600$ wiederholt wird. Drittens
durchläuft der $\chi_{21}$-Gap einen Vorzeichenwechsel zwischen $N=200$ und $N=600$
($+0.282\to-0.004$), was veranschaulicht, dass selbst der Galerkin-Operator
unterhalb seiner Konvergenzskala nicht informativ ist. Wir katalogisieren die
zehn Charaktere anhand zweier \emph{konstanter Filtermerkmale} (Parität und Wurzelzahl,
beide im gesamten Sample gleich $+1$) und zweier \emph{informativer Achsen}
(archimedisches Gewichtsverhältnis und Gap-Signatur), zeigen die systematische
Siegel-Walfisz-Kompression bei
$N=600$ und positionieren das Ergebnis als präzise Diagnose, warum
numerische Galerkin-Routen für sich genommen
$C^{(2)}_{\chifn}$ nicht auflösen können. Eine echte prädiktive Formulierung würde eine
asymptotische Paley-Wiener-Kontrolle der Sektor-Grundzustände erfordern,
die die Galerkin-Trunkierung umgeht; wir schließen mit einem Fahrplan hin zu
einem solchen branch-lokalen Hilbert-P\'olya-Rahmen (branch-local Hilbert-P\'olya framework).
\end{abstract}
\vspace{1em}
]
\clearpage

\section{Einführung}
\label{sec:intro}

\subsection{Motivation}
\label{sec:intro-motivation}

Der aktuelle Reduktionsansatz über gerade Dominanz zur Riemannschen
Vermutung \cite{RH_II} stellt den Sektor gerader Parität der Weilschen
quadratischen Form auf $L^2[-L,L]$, $L=\log\lambda$, als Route zu einer
uniformen Dominanz über den ungeraden Sektor dar, mit der Zielrate
$\mathrm{gap}(\lambda)\gtrsim\sqrt{\lambda}$. Die Konstruktion wird dort
als bedingte Reduktion und diagnostische Route gerahmt, nicht als
abgeschlossener unbedingter Beweis; sie ist variationell und verwendet
lediglich die drei strukturellen Bestandteile von $\zeta(s)$: das
Euler-Produkt, den Gamma-Faktor und eine einzige Funktionalgleichung.

Eine natürliche und zentrale Frage ist, ob sich dieselbe Architektur
auf getwistete $L$-Funktionen $L(s,\chi)$ für nicht-triviale
Dirichlet-Charaktere $\chi$ übertragen lässt, um so die Verallgemeinerte
Riemannsche Vermutung (Generalized Riemann Hypothesis, GRH) charakterweise zu adressieren. Im Meta-Programm
von \cite{MetaGalerkinBridge} wird die RH-Typ-Frage für eine Zeta-Familie
$X$ durch eine Dreikomponenten-Zerlegung diagnostiziert:
\begin{equation}
\label{eq:three-component}
  \mathrm{RH}_X
    \;=\;
  \mathrm{GBC}_X \;+\; C^{(2)}_X \;+\; \mathrm{Hurwitz}_X,
\end{equation}
wobei $\mathrm{GBC}_X$ die Galerkin-Brücke zu einem selbstadjungierten
Operator erfasst, $\mathrm{Hurwitz}_X$ den Realität-der-Nullstellen-Transfer, und
$C^{(2)}_X$ die Operator-Identifikation im Connes-Stil, welche das Herzstück des
Hilbert-P\'olya-Programms bildet. Für den trivialen Charakter isoliert
\cite{RH_II} die entsprechende Route über gerade Dominanz für
$C^{(2)}_{X=\zeta}$; der Abschluss wird dort als bedingte Reduktion und
nicht als unbedingtes Theorem geführt. Für nicht-triviale $\chi$ ist
$C^{(2)}_{\chifn}$ die harte Residualkomponente. Die vorliegende Arbeit
kartiert, was Galerkin-numerische Techniken über dieses Residuum aussagen
können und, noch wichtiger, was sie nicht aussagen können.

\subsection{Was diese Arbeit leistet}
\label{sec:intro-contribution}

Wir katalogisieren den Galerkin-Sektor-Gap
$\gap_{\chifn}(\lambda)
  :=\lambda_{\min}(Q^-_{\chifn})-\lambda_{\min}(Q^+_{\chifn})$
bei $\lambda=20000$ für zehn Kronecker-reelle Charaktere
$\chi_D$, $D\in\{5,8,12,13,17,21,24,29,33,60\}$,
über drei Galerkin-Auflösungen $N\in\{200,400,600\}$ hinweg. Für jeden
Charakter berechnen wir beide Sektor-Grundzustände separat
$\phi^\pm_{\chifn}$, die primitive Weil-Kern-Prädiktion
\begin{equation}
\label{eq:predictor}
  S_{\chifn}(\lambda)
    = -2\sum_{\substack{p\nmid q\\m\log p\leq 2L}}
      \frac{\chifn(p)^m\log p}{p^{m/2}}\;\Deltachi(m\log p),
\end{equation}
mit
$\Deltachi(t)
  =(\phi^+_{\chifn}\!*\phi^+_{\chifn})(t)
    +(\phi^-_{\chifn}\!*\phi^-_{\chifn})(t)$,
den archimedischen Beitrag
$\archdiff_{\chifn}
  :=\langle\phi^-,K^{\mathrm{arch}}_-\phi^-\rangle
    -\langle\phi^+,K^{\mathrm{arch}}_+\phi^+\rangle$,
und den vollen Galerkin-Gap $\gap_{\mathrm{gal}}^{(N)}$. Wir positionieren die
vier Größen in Relation zueinander, über alle drei Galerkin-Auflösungen
hinweg, sowie als Referenz in Relation zu den empirischen Gaps (bei höchstem $N$).

\subsection{Was diese Arbeit aufzeigt}
\label{sec:intro-findings}

Unser Beitrag ist bewusst diagnostischer Natur. Die wichtigsten Erkenntnisse sind:

\begin{enumerate}
\item[(F1)] \textbf{Falsifizierung der Reduktion führender Ordnung.}
  Unter der Approximation
  $\phi^+_{\chifn}\approx\phi^-_{\chifn}\approx\phi_\infty$, die in
  \cite{RH_II} für den trivialen Charakter verwendet wird, reduziert sich der Primzahl-Prädiktor
  auf ein universelles Autokorrelations-Integral. Ein systematischer Test
  von 16 Fenstermodellen, einschließlich des empirischen Einzelsektor-Grundzustands,
  erreicht bestenfalls eine Vorzeichengenauigkeit von $6/10$ und ein $R^2=0.15$.
  Die Reduktion führender Ordnung ist für nicht-triviale $\chi$ falsch.
\item[(F2)] \textbf{Matrixzerlegungs-Tautologie.}
  Wenn \emph{beide} Sektor-Grundzustände aus derselben
  Galerkin-Diagonalisierung entnommen werden, ist die Identität
  $\gap_{\mathrm{gal}}^{(N)} = S_{\chifn}+\archdiff_{\chifn}$
  keine Vorhersage, sondern eine Folge der Linearität der
  Erwartung auf $W=W_{\mathrm{arch}}+W_{\mathrm{prime}}$. Eine scheinbare
  $9/9$ Vorzeichengenauigkeit bei $N=200$ (abzüglich des bekannten $N$-Oszillators
  $\chi_{21}$) löst sich bei der Validierung auf Serverebene für $N=600$ auf: Der
  Prädiktor $S_{\chifn}+\archdiff_{\chifn}$ fällt auf $8/10$ aufgrund
  zweier Vorzeichenfehler: $\chi_{21}$ ($N$-Oszillator,
  vgl. (F3) und \S\ref{sec:chi21}) sowie $\chi_{29}$ (numerisches
  Interpolationsartefakt an einem Gap nahe null).
\item[(F3)] \textbf{$\chi_{21}$ als ein $N$-Oszillator.}
  Bei $N=200$ hat man $\gap_{\mathrm{gal}}=+0.282$; bei $N=600$ hat man
  $\gap_{\mathrm{gal}}=-0.004$. Der Vorzeichenwechsel ist kein Formelfehler,
  sondern ein Galerkin-Trunkierungseffekt für zusammengesetzte Führer, der sich erst
  oberhalb von $N\sim 500$ stabilisiert. Dies zeigt einerseits, dass der Galerkin-Operator
  selbst nicht-trivial $N$-abhängig ist, und erklärt andererseits, warum ein
  $9/9$ bei $N=200$ nach Ausschluss von $\chi_{21}$ nicht überraschend war:
  Der ausgeschlossene Charakter war genau jener, dessen $N$-Abhängigkeit noch nicht
  in den Test integriert war.
\item[(F4)] \textbf{Residualer Inhalt: Die Weil-Kern-Architektur ist
  rigoros.} Die Sektor-Kern-Zerlegung
  (Theorem~\ref{thm:kernel}), die anti-diagonale Sektordifferenz
  (Theorem~\ref{thm:sector-diff}) und die spektral-domänische Form des
  Gaps über die $L'/L$-Funktion bleiben als strukturelle Aussagen gültig,
  unabhängig von der Reduktion führender Ordnung. Wir legen diese in
  Abschnitt~\ref{sec:setup} dar und verwenden sie in Abschnitt~\ref{sec:atlas},
  um die zehn Charaktere entlang vier orthogonaler Achsen zu klassifizieren.
\end{enumerate}

Zusammengenommen liefern diese Erkenntnisse eine präzise Diagnose: Die
$C^{(2)}_{\chifn}$-Komponente von \eqref{eq:three-component} ist für
nicht-triviale $\chi$ einer Galerkin-numerischen Schließung
nicht zugänglich, da das relevante Signal in der Differenz
$\phi^+_{\chifn}-\phi^-_{\chifn}$ lebt; dies ist eine Feinstruktur, die in
jeder natürlichen Approximation führender Ordnung verloren geht, und überdies ist der volle
Galerkin-Operator relativ zu seinen eigenen Eigenvektoren tautologisch.
Eine echte prädiktive Formulierung muss
$\Deltachi$ analytisch identifizieren, z.\,B.\ durch Paley-Wiener-Asymptotik
\cite{PaleyWiener}, angepasst aus \cite{RH_II} auf das
charakter-getwistete Setting.

\subsection{Organisation}

Abschnitt~\ref{sec:setup} fasst die Weilsche quadratische Form und die
Sektor-Kern-Zerlegung zusammen. Abschnitt~\ref{sec:formula} präsentiert die
Gap-Formel und kontrastiert die reduzierte und unreduzierte Version.
Abschnitt~\ref{sec:validation} präsentiert die numerische Validierung bei
$N\in\{200,400,600\}$. Abschnitt~\ref{sec:atlas} katalogisiert die zehn
Charaktere entlang vier Achsen. Abschnitt~\ref{sec:tautology} formuliert
und beweist die Tautologie. Abschnitt~\ref{sec:chi21} behandelt die
$\chi_{21}$ $N$-Oszillation im Detail. Abschnitt~\ref{sec:discussion}
diskutiert Implikationen für \eqref{eq:three-component}. Abschnitt~\ref{sec:outlook}
skizziert den Paley-Wiener-Fahrplan und das branch-lokale
Hilbert-P\'olya-Programm als den angezeigten nächsten Schritt.

\section{Grundlagen: Weilsche quadratische Form und Sektor-Kerne}
\label{sec:setup}

Wir fassen den durchgehend verwendeten strukturellen Rahmen zusammen; die vollständigen
Herleitungen sind in \cite{AnalyticKernel} gegeben.

\subsection{\texorpdfstring{Dirichlet $L$-Funktionen}{Dirichlet L-Funktionen} und die Weil-Formel}

Sei $\chifn$ ein primitiver Dirichlet-Charakter mod $q$ mit
$\chifn(-1)=+1$ (gerade). Setze $a_{\chifn}=0$ und definiere die vervollständigte
$L$-Funktion
\[
  \Lambda(s,\chifn)
    =\Bigl(\frac{q}{\pi}\Bigr)^{s/2}\Gamma\!\Bigl(\frac{s}{2}\Bigr)
      L(s,\chifn),
\]
welche die Funktionalgleichung
$\Lambda(s,\chifn)=\epsilon_{\chifn}\Lambda(1-s,\bar\chifn)$ mit
$|\epsilon_{\chifn}|=1$ erfüllt. Für einen reellen (d.\,h.\ Kronecker) Charakter
ist $\bar\chifn=\chifn$ und $\epsilon_{\chifn}=+1$; wir arbeiten in diesem Atlas ausschließlich
mit den zehn reellen Charakteren $\chifn_D$ mit $D\in\{5,8,12,13,17,21,
24,29,33,60\}$. Die nicht-trivialen Nullstellen
$\rho=\halfline+i\gamma$ werden durchgehend auf der kritischen Gerade angenommen
(GRH für die getesteten Führer). Die Datenbankseiten von \cite{LMFDB}
liefern berechnete niedrig liegende Nullstellen auf der kritischen
Geraden für die entsprechenden Dirichlet-$L$-Funktionen. Wir verwenden
diese Einträge nur als endliche numerische Konsistenzdaten; sie sind
kein Zertifikat dafür, dass alle Nullstellen unterhalb einer bestimmten
Höhe auf der kritischen Geraden verifiziert wären. Ein theoretischer
Beweis von GRH für $L(s,\chifn_D)$ für irgendeinen spezifischen $D$ in
der Stichprobe steht selbstverständlich aus.

Für eine gerade Testfunktion $h\in C_c^\infty(\mathbb{R})$ mit kompaktem Träger
lautet die explizite Weilsche Formel (Weil explicit formula)
\cite{Weil1952,IwaniecKowalski}:
\begin{equation}
\label{eq:weil}
  \sum_{\gamma}h(\gamma)
    = W^{\chifn}_\infty[h]
    - 2\sum_{\substack{p\nmid q\\m\geq 1}}
      \frac{\chifn(p)^m\log p}{p^{m/2}}\,\hat h(m\log p),
\end{equation}
wobei die Summe auf der linken Seite über alle Imaginärteile der nicht-trivialen
Nullstellen von $L(s,\chifn)$ läuft und das archimedische Gewicht lautet:
\[
  W^{\chifn}_\infty[h]
    = \frac{1}{2\pi}\!\int_{-\infty}^{\infty}\! h(t)
      \Bigl[\log\!\frac{q}{\pi} + \Re\,\psi\!\Bigl(\tfrac14+\tfrac{it}{2}\Bigr)\Bigr]\,dt.
\]

\subsection{\texorpdfstring{Quadratische Form auf $H_L$}{Quadratische Form auf HL}}

Fixiere $\lambda>1$ und setze $L:=\log\lambda$, $H_L:=L^2[-L,L]$. Für eine
reelle Testfunktion $f\in C_c^\infty([-L,L])$, setze
$h(t):=2\pi|\hat f(t)|^2$. Dann ist $h$ nicht-negativ und gerade, und
\eqref{eq:weil} induziert die \emph{Weilsche quadratische Form}
\begin{equation}
\label{eq:weil-qf}
\begin{aligned}
  Q_{\chifn}[f]
    &:=\sum_{\gamma}2\pi|\hat f(\gamma)|^2 \\
    &=W^{\chifn}_\infty[2\pi|\hat f|^2] \\
    &\quad -2\sum_{p,m}\frac{\chifn(p)^m\log p}{p^{m/2}}
      \int f(y)f(y-m\log p)\,dy.
\end{aligned}
\end{equation}
$Q_{\chifn}$ lässt ein translationsinvariantes Symbol zu: Durch Definition der
primitiven Faltungsdistribution (primitive convolution distribution)
\begin{equation}
\label{eq:kappa}
\begin{aligned}
  \kappa_{\chifn}(u)
    &:=-\sum_{p\nmid q,\,m\geq 1}
      \frac{\chifn(p)^m\log p}{p^{m/2}} \\
    &\quad\times
      \bigl[\delta(u-m\log p)+\delta(u+m\log p)\bigr],
\end{aligned}
\end{equation}
und des archimedischen Kerns
$K^{\mathrm{arch}}_{\chifn}(u):=\mathcal{F}^{-1}[\rho_{\chifn}](u)$
mit $\rho_{\chifn}(t)=\log(q/\pi)+\Re\,\psi(1/4+it/2)$ erhält man
den Kern
$K_{\chifn}(x,y)=\kappa_{\chifn}(x-y)+K^{\mathrm{arch}}_{\chifn}(x-y)$,
der $Q_{\chifn}[f]=\iint K_{\chifn}(x,y)f(x)f(y)\,dx\,dy$ auf
$C_c^\infty([-L,L])$ erfüllt.

\subsection{Sektor-Zerlegung}

Sei $\mathcal{P}:H_L\to H_L$, $(\mathcal{P}f)(x)=f(-x)$, die
Spiegelung. $\mathcal{P}$ ist unitär selbstadjungiert mit $\mathcal{P}^2
=I$; die Spektralprojektionen $P^\pm=(I\pm\mathcal{P})/2$ spalten
$H_L=H_L^+\oplus H_L^-$ in gerade/ungerade Sektoren. Standard-Fourierbasen
sind
\[
\begin{aligned}
  e_0^+&=\tfrac{1}{\sqrt{2L}}, &
  e_n^+(x)&=\tfrac{1}{\sqrt L}\cos(n\pi x/L),\\
  e_n^-(x)&=\tfrac{1}{\sqrt L}\sin(n\pi x/L).
\end{aligned}
\]
Die Galerkin-Trunkierung bei der Basisgröße $N$ verwendet die Moden $0\leq n\leq N$
in jedem Sektor.

\subsection{Sektor-Kerne und die anti-diagonale Differenz}

Da $K_{\chifn}(x,y)=K_{\chifn}(-x,-y)$, nehmen die sektor-projizierten
Kerne folgende Form an:
\begin{theorem}[Sektor-Kerne; {\cite{AnalyticKernel}, Thm.~3.2}]
\label{thm:kernel}
Auf $H_L^\pm$ gilt jeweils:
\[
\begin{aligned}
  K_{\chifn}^{\pm}(x,y)
    &=\tfrac{1}{2}\bigl[\kappa_{\chifn}(x-y)
       \pm\kappa_{\chifn}(x+y)\bigr]\\
    &\quad+\tfrac{1}{2}\bigl[K^{\mathrm{arch}}_{\chifn}(x-y)
       \pm K^{\mathrm{arch}}_{\chifn}(x+y)\bigr].
\end{aligned}
\]
\end{theorem}

\begin{proof}[Beweisskizze]
Die Reflexion $\mathcal{P}f(x):=f(-x)$ vertauscht mit Translation
$T_a$ via $\mathcal{P}T_a=T_{-a}\mathcal{P}$, und der Kern
$K_\chi(x,y)=\kappa_\chi(x-y)+K^{\mathrm{arch}}_\chi(x-y)$ ist
translationsinvariant. Symmetrisierung unter $\mathcal{P}$:
$P^\pm K_\chi P^\pm = (K_\chi \pm \mathcal{P}K_\chi\mathcal{P})/2$,
wobei $\mathcal{P}K_\chi\mathcal{P}(x,y)=K_\chi(-x,-y)=K_\chi(x,y)$
wegen der Parität von $\kappa_\chi$ und $K^{\mathrm{arch}}_\chi$
(beide gerade auf $\mathbb{R}$). Die Off-Diagonal-Symmetrisierung
$P^\pm K_\chi (1-P^\pm)$ verschwindet auf dem entsprechenden Sektor
und liefert den dargestellten Ausdruck. Die Wirkung auf den
Test-Funktionen-Raum $C_c^\infty([-L,L])$ ist wohldefiniert, weil
$\kappa_\chi$ als temperierte Distribution auf $\{u=\pm m\log p\}$
getragen ist, alles in $[-2L,2L]$; Paarung mit $f\in C_c^\infty([-L,L])$
liefert eine endliche Summe über $(p,m)$ mit $m\log p\leq 2L$. Volle
Rechendetails inklusive der Verifikation, dass die Symmetrisierung
Selbstadjungiertheit erhält, finden sich in \cite{AnalyticKernel},
\S3.2.
\end{proof}

Projizierte Eigenwertprobleme
$K_{\chifn}^{\pm}\phi_{\chifn}^{\pm}=\lambda_1(\chifn,\pm)\phi_{\chifn}^{\pm}$
auf $H_L^\pm$ definieren die Sektor-Grundzustände $\phi_{\chifn}^\pm$,
die im Folgenden eine zentrale Rolle spielen.

\begin{theorem}[Anti-diagonale Sektor-Differenz; {\cite{AnalyticKernel}, Thm.~4.1}]
\label{thm:sector-diff}
Für einen geraden primitiven Dirichlet-Charakter $\chifn$ (d.\,h.\
$\chifn(-1)=+1$, wie im gesamten Atlas) gilt
\[
  K_{\chifn}^-(x,y)-K_{\chifn}^+(x,y)
    =-\kappa_{\chifn}(x+y)-K^{\mathrm{arch}}_{\chifn}(x+y).
\]
Der primitive Teil ist auf den Anti-Diagonalen $x+y=\pm m\log p$ getragen
und enthält die volle Charakter-Abhängigkeit; der archimedische Teil
hängt von $\chifn$ nur durch die Konstante $\log(q/\pi)$ ab (im
gerade-Charakter-Setting; im ungeraden Fall käme ein
$\Re\,\psi(3/4+it/2)$-Term hinzu, der außerhalb des Geltungsbereichs
dieses Atlasses liegt).
\end{theorem}

\begin{proof}[Beweisskizze]
Aus Theorem~\ref{thm:kernel} folgt durch Subtraktion
$K_\chi^-(x,y)-K_\chi^+(x,y)$: Alle $(x-y)$-Argumente heben sich auf,
die $(x+y)$-Argumente verdoppeln sich mit Vorzeichen $-$:
$K_\chi^- - K_\chi^+ = -\kappa_\chi(x+y) - K^{\mathrm{arch}}_\chi(x+y)$.
Zur Charakterabhängigkeit: $\kappa_\chi$ hängt von $\chifn$ über jeden
$\chifn(p)^m\log p / p^{m/2}$-Koeffizienten ab (\eqref{eq:kappa}),
während $K^{\mathrm{arch}}_\chi=\mathcal{F}^{-1}[\rho_\chi]$ mit
$\rho_\chi(t)=\log(q/\pi)+\Re\,\psi(1/4+it/2)$ nur über die Konstante
$\log(q/\pi)$ von $\chifn$ abhängt (der $\psi$-Term ist im geraden
Setting charakter-unabhängig). Vollständige Herleitung:
\cite{AnalyticKernel}, \S4.1.
\end{proof}

\subsection{Das Gap-Funktional}

Sei $\lambda_1^{\min}(Q_{\chifn}^\pm)$ der kleinste Eigenwert
von $Q_{\chifn}$ eingeschränkt auf $H_L^\pm$, und schreibe
\begin{equation}
\label{eq:gap-def}
\begin{aligned}
  \gap_{\chifn}(\lambda)
    &:=\lambda_1^{\min}(Q_{\chifn}^-)
      -\lambda_1^{\min}(Q_{\chifn}^+)\\
    &=\langle\phi_{\chifn}^-|K_{\chifn}^-|\phi_{\chifn}^-\rangle
      -\langle\phi_{\chifn}^+|K_{\chifn}^+|\phi_{\chifn}^+\rangle.
\end{aligned}
\end{equation}
Für den trivialen Charakter (klassisches Riemann-Setting) zielt
\cite{RH_II} auf eine $\sqrt{\lambda}$-Sektortrennung als Mechanismus
gerader Dominanz; die aktuell zitierte Fassung rahmt dies als bedingte
Reduktion und nicht als abgeschätzten unbedingten Abschluss. Für
nicht-triviale $\chifn$ lautet die empirische Beobachtung, dass
$|\gap_{\chifn}|=O(1)$ bei zugänglichem $\lambda$, mit
einem \emph{charakter-spezifischen} Vorzeichen; der Atlas in
Abschnitt~\ref{sec:validation} ist das quantitative Abbild dieser
Beobachtung.

\section{Die Gap-Formel: reduziert vs.\ unreduziert}
\label{sec:formula}

\subsection{Die Reduktion führender Ordnung}
\label{sec:formula-reduced}

Ausgehend von \eqref{eq:gap-def} und
Theorem~\ref{thm:sector-diff} lässt die Lücke folgende Darstellung zu:
\begin{equation}
\label{eq:gap-matrix}
\begin{aligned}
  \gap_{\chifn}(\lambda)
    &= \langle\phi_{\chifn}^-|K_{\chifn}^--K_{\chifn}^+|\phi_{\chifn}^-\rangle\\
    &\quad + \langle\phi_{\chifn}^-|K_{\chifn}^+|\phi_{\chifn}^-\rangle
      - \langle\phi_{\chifn}^+|K_{\chifn}^+|\phi_{\chifn}^+\rangle.
\end{aligned}
\end{equation}
Approximiert man beide Grundzustände durch eine gemeinsame Funktion
$\phi_{\chifn}^\pm\approx\phi_\infty$, verschwindet die zweite Klammer und
die erste vereinfacht sich nach Theorem~\ref{thm:sector-diff} zu
\begin{equation}
\label{eq:gap-reduced}
\begin{aligned}
  \gap_{\chifn}(\lambda)
    &\overset{?}{\approx}
      -\iint \phi_\infty(x)\bigl[\kappa_{\chifn}(x+y)
      +K^{\mathrm{arch}}_{\chifn}(x+y)\bigr]\\
    &\quad\times\phi_\infty(y)\,dx\,dy.
\end{aligned}
\end{equation}
Das Entfalten von $\kappa_{\chifn}$ über \eqref{eq:kappa} und die Nutzung der Geradheit
(evenness) von $\phi_\infty$ (gerader Sektor) liefert die reduzierte Gap-Formel
\begin{equation}
\label{eq:gap-reduced-explicit}
\begin{aligned}
  \gap_{\chifn}(\lambda)
    &\overset{?}{\approx}
      2\sum_{\substack{p\nmid q\\m\log p\leq 2L}}
        \frac{\chifn(p)^m\log p}{p^{m/2}}\\
    &\quad\times\Phi_\infty(m\log p)
      \;+\;\text{arch.\ term},
\end{aligned}
\end{equation}
mit $\Phi_\infty(t):=(\phi_\infty\ast\phi_\infty)(t)$
(vgl.\ \cite{AnalyticKernel}, Cor.~4.4 und \cite{AnalyticGroundstate},
Cor.~6.1). Dies ist die Formel, die man natürlicherweise vom
trivialen Charakter-Fall übertragen würde, wo der gemeinsame Grundzustand $\phi_\infty$
aus der Paley-Wiener-Analyse von \cite{RH_II} zugänglich ist.

\subsection{Empirischer Status der reduzierten Formel}
\label{sec:formula-reduced-empirical}

In Session 6 des vorliegenden Programms wurde \eqref{eq:gap-reduced-explicit}
mit sechzehn Kandidatenfenstern $\phi_\infty$ getestet: zentrierte Gaußfunktionen der
Breite $\sigma\in\{0.1,0.3,0.5,1.0,2.0\}$, verschobene Gaußfunktionen, Rechtecke
mit Cut-off $T$, $\cos^2$-Fenster auf $[-T,T]$, die rohe
$\chifn$-gewichtete Primsumme, und---am entscheidendsten---der \emph{empirische}
Einzelsektor-Grundzustand (single-sector ground state)
$\phi_\infty:=\phi_{\chifn}^+$ extrahiert aus der $N=200$ Galerkin-Diagonalisierung.

Keiner der sechzehn Kandidaten erreichte das Erfolgskriterium von
$\geq\!9/10$ Vorzeichenübereinstimmungen mit $R^2\geq 0.8$. Das beste
Gaußfenster führender Ordnung erreicht $6/10$ mit $R^2=0.15$; der
empirische Einzelsektor-Grundzustand erreicht $4/10$ mit
$R^2=0.02$. Details des 16-Fenster-Scans sind in
\cite{Woche4Validation} festgehalten. Die entscheidende Beobachtung ist, dass für alle
zehn Charaktere die empirische Einzelsektor-Autokorrelation
$\Phi_\infty$ nahezu identisch ist (numerisch innerhalb von $10^{-3}$):
Die Reduktion $\phi_{\chifn}^+\approx\phi_{\chifn}^-\approx
\phi_\infty$ ist somit nicht nur quantitativ ungenau, sondern
löscht das charakter-spezifische Signal.

\begin{remark}[Wo das Charakter-Signal lebt]
\label{rem:signal-locus}
Die auf den Sektor eingeschränkten Grundzustände sind in führender Ordnung dieselbe
Funktion: Sie liegen beide nah am Paley-Wiener Grundzustand des
diagonalen Kerns $\kappa_{\chifn}(x-y)+K^{\mathrm{arch}}_{\chifn}
(x-y)$. Charakter-spezifische Informationen werden in der Differenz
$\phi_{\chifn}^+-\phi_{\chifn}^-$ gespeichert, was der Feinstruktur
entspricht, die durch den anti-diagonalen Term $\kappa_{\chifn}(x+y)+
K^{\mathrm{arch}}_{\chifn}(x+y)$ von Theorem~\ref{thm:sector-diff} erzeugt wird.
Jede analytische Reduktion, die $\phi_{\chifn}^+$ und
$\phi_{\chifn}^-$ in eine gemeinsame Funktion absorbiert, verliert exakt dieses Signal.
\end{remark}

\subsection{Die unreduzierte Formel}
\label{sec:formula-unreduced}

Wir behalten beide Grundzustände separat bei. Definiere die Sektordifferenz-Autokorrelation
\begin{equation}
\label{eq:Delta-def}
  \Deltachi(t)
    := (\phi^+_{\chifn}\!*\phi^+_{\chifn})(t)
      + (\phi^-_{\chifn}\!*\phi^-_{\chifn})(t),
\end{equation}
\emph{Parität-Vorzeichen-Hinweis.} Die numerische Implementierung in unseren Skripten
(Anhang~\ref{app:sign}) berechnet die Autokorrelation
$R_{\phi^\pm}(t):=\int\phi^\pm(s)\phi^\pm(s-t)\,ds$ mittels
\texttt{np.correlate} und bildet dann
$\Deltachi^{\mathrm{num}}(t):=R_{\phi^+}(t)-R_{\phi^-}(t)$.
Aus der Paritätsidentität $R_{\phi^+}=(\phi^+{*}\phi^+)_{\mathrm{math}}$
(gerader Sektor) und $R_{\phi^-}=-(\phi^-{*}\phi^-)_{\mathrm{math}}$
(ungerader Sektor) folgt
\[
  \Deltachi^{\mathrm{num}}(t)
    =(\phi^+{*}\phi^+)_{\mathrm{math}}(t)+(\phi^-{*}\phi^-)_{\mathrm{math}}(t),
\]
was genau der rechten Seite von~\eqref{eq:Delta-def} entspricht. Damit stimmt
$\Deltachi$ wie in~\eqref{eq:Delta-def} definiert identisch mit der
skript-berechneten Größe $\Deltachi^{\mathrm{num}}$ überein, und der
Koeffizient~$-2$ in~\eqref{eq:S-def} ist der eindeutige Wert, der
die Gap-Identität (Satz~\ref{thm:closure}) erfüllt. Siehe
Anhang~\ref{app:sign} für die vollständige Herleitung mit allen sechs
relevanten Vorzeichenquellen.

und die archimedische Sektordifferenz
\begin{equation}
\label{eq:arch-diff-def}
  \archdiff_{\chifn}
    := \langle\phi^-_{\chifn},K^{\mathrm{arch}}_-\phi^-_{\chifn}\rangle
     - \langle\phi^+_{\chifn},K^{\mathrm{arch}}_+\phi^+_{\chifn}\rangle,
\end{equation}
wobei $K^{\mathrm{arch}}_\pm$ die sektor-eingeschränkten archimedischen
Operatoren sind (Theorem~\ref{thm:kernel}). Einsetzen in
\eqref{eq:gap-matrix} und Nutzung von Theorem~\ref{thm:sector-diff} liefert
die untenstehende unreduzierte Formel.

\begin{theorem}[Unreduzierte Gap-Formel]
\label{thm:closure}
Sei $\chifn$ ein primitiver reeller Charakter mit Führer $q$ und
$\chifn(-1)=+1$, und sei $\lambda>1$, $L=\log\lambda$. Seien
$\phi^\pm_{\chifn}\in H_L^\pm$ die Sektor-Grundzustände von
$K_{\chifn}^\pm$, jeweils $L^2$-normalisiert. Dann gilt
\begin{equation}
\label{eq:gap-unreduced}
  \gap_{\chifn}(\lambda)
    = S_{\chifn}(\lambda) + \archdiff_{\chifn},
\end{equation}
mit
\begin{equation}
\label{eq:S-def}
  S_{\chifn}(\lambda)
    = -2\sum_{\substack{p\nmid q\\m\log p\leq 2L}}
      \frac{\chifn(p)^m\log p}{p^{m/2}}\,\Deltachi(m\log p).
\end{equation}
\end{theorem}

Eine Herleitung findet sich in \cite{AnalyticKernel}, Bemerkung~4.5 (mit
der Vorzeichenkonvention abgeglichen mit \texttt{build\_W}; siehe
Anhang~\ref{app:sign}). Auf dem Niveau der \emph{Galerkin-Trunkierung}
bei Basisgröße $N$ gilt Gleichung
\eqref{eq:gap-unreduced} ebenfalls exakt:
\begin{equation}
\label{eq:gap-galerkin}
  \gap_{\mathrm{gal}}^{(N)}(\chifn)
    = S_{\chifn}^{(N)}(\lambda) + \archdiff_{\chifn}^{(N)},
\end{equation}
vorausgesetzt $\phi^\pm_{\chifn}$ werden als die Grundzustands-Eigenvektoren
der Galerkin-trunkierten Operatoren $K_{\chifn}^{\pm, N}$
bei Basisgröße $N$ genommen und $S^{(N)}$, $\archdiff^{(N)}$ werden
aus diesen endlichen-$N$-Daten berechnet.

\subsection{Zwei numerische Realisierungen von \texorpdfstring{$S_{\chifn}^{(N)}$}{S}}
\label{sec:formula-two-realizations}

Ein subtiler Punkt, der alles Folgende antreibt, ist, dass
\eqref{eq:S-def} \emph{zwei} verschiedene numerische
Implementierungen aus Galerkin-Daten zulässt, die bis auf $10^{-3}$ übereinstimmen,
jedoch nicht identisch sind.
\begin{itemize}
  \item[(R1)] \emph{Räumliche Autokorrelation (Spatial autocorrelation).} Werte
    $\phi_{\chifn}^\pm$ auf einem Raster (grid) aus, berechne $(\phi^\pm
    *\phi^\pm)(t)$ mit einer numerischen Korrelation, interpoliere auf
    die Primzahlabzissen $t=m\log p$, und summiere.
  \item[(R2)] \emph{Matrix-Zerlegung.} Zerlege die Galerkin-Matrix
    $W^\pm=W^\pm_{\mathrm{arch}}+W^\pm_{\mathrm{prime}}$ und
    bilde
    $\mathrm{prime\_diff}_{\mathrm{matrix}}
      :=\langle\phi^-,W^-_{\mathrm{prime}}\phi^-\rangle
      -\langle\phi^+,W^+_{\mathrm{prime}}\phi^+\rangle$.
\end{itemize}
Beide Größen schätzen $S_{\chifn}^{(N)}$ ab; (R1) wird über
interpolierte Autokorrelation erhalten, (R2) über ein exaktes Matrix-Skalarprodukt.
Ihre Diskrepanz liegt auf dem Niveau des Gleitkomma- und
Interpolationsfehlers in der Größenordnung $10^{-3}$. Für Charaktere mit
$|\gap_{\chifn}|\gtrsim 10^{-2}$ ist dies vernachlässigbar. Für die schwach
signalisierten Charaktere $\chifn_{21}$, $\chifn_{29}$, $\chifn_{60}$,
mit $|\gap|\lesssim 10^{-2}$, ist es in der gleichen Größenordnung wie das
Signal. Abschnitt~\ref{sec:tautology} widmet sich den Konsequenzen.

\section{Numerische Validierung}
\label{sec:validation}

\subsection{Berechnungsaufbau}
\label{sec:validation-setup}

Für jeden der zehn reellen Charaktere $\chifn_D$,
$D\in\{5,8,12,13,17,21,24,29,33,60\}$, konstruieren wir die Galerkin-Matrizen
$W^\pm_{\chifn}$ in den Fourier-Basen von $H_L^\pm$ beim Cut-off
$\lambda=20000$ ($L=\log\lambda\approx 9.90$). Die Assemblierung der Matrix (matrix assembly) verwendet die
Koeffizienten der expliziten Formel $\chifn(p^m)\log p/p^{m/2}$ für alle
Primzahlpotenzen mit $m\log p\leq 2L$ und den archimedischen Operator
$K^{\mathrm{arch}}_{\chifn}$, berechnet mittels Fast Fourier Transform von
$\rho_{\chifn}$.

Drei Galerkin-Auflösungen werden verwendet: $N=200$ lokal auf einem Desktop,
$N=400$ lokal, und $N=600$ auf einer 8\,GB Hetzner CCX13 Cloud-Instanz
(Gesamt-Wall-Clock 1.95 Stunden über alle zehn Charaktere). Der
kleinste Eigenwert jedes $W^\pm_{\chifn,N}$ wird durch eine
Lanczos-Inverse-Iteration mit Toleranz $10^{-10}$ ermittelt, was den
Sektor-Grundzustand $\phi^{\pm,(N)}_{\chifn}$ und den Galerkin-Gap
$\gap_{\mathrm{gal}}^{(N)}(\chifn)
:=\lambda_1^{\min}(W^-_{\chifn,N})-\lambda_1^{\min}(W^+_{\chifn,N})$ liefert.

Für jeden Charakter werden die Größen $S_{\chifn}^{(N)}$ (via
Räumlicher-Autokorrelations-Realisierung (R1)) und
$\archdiff_{\chifn}^{(N)}$ aus den Grundzuständen ausgewertet, und
die Summe $S_{\chifn}^{(N)}+\archdiff_{\chifn}^{(N)}$ wird mit
$\gap_{\mathrm{gal}}^{(N)}(\chifn)$ sowie mit dem Galerkin-Referenz-Gap (highest-$N$)
$\gap_{\mathrm{gal}}^{(\mathrm{ref})}(\chifn)$ verglichen. Die Skripte
\texttt{arch\_term\_analysis.py} und
\texttt{server\_n600\_full\_analysis.py} speichern die Rohausgaben in
\texttt{ARCH\_TERM\_ANALYSIS.json} und
\texttt{ARCH\_TERM\_N600\_SERVER.json}; Tabellen zur Nachbearbeitung
finden sich in \texttt{ARCH\_TERM\_ANALYSIS.md} und
\texttt{ARCH\_TERM\_N600\_ANALYSIS.md}.

\subsection{Die zentrale Ergebnistabelle}
\label{sec:validation-table}

Tabelle~\ref{tab:validation} fasst den Vergleich bei $\lambda=
20000$ für alle zehn Charaktere zusammen. Für jeden Charakter listen wir den
Galerkin-Referenz-Gap $\gap_{\mathrm{gal}}^{(\mathrm{ref})}$ (mit der Auflösung
$N_{\mathrm{src}}$, aus der er stammt), den Galerkin-Gap bei
$N=200$ und $N=600$, sowie den Prädiktor $S+\archdiff$ bei beiden
Auflösungen auf. Zwei Beobachtungen verankern die gesamte Analyse.

Erstens stimmen bei $N=600$ die Einträge mit der Bezeichnung „$\gap_{\mathrm{gal}}^{(\mathrm{ref})}$“ und
„$\gap_{\mathrm{gal}}^{(600)}$“ immer dann identisch überein, wenn die
empirische Referenz selbst bei $N=600$ berechnet wurde (die sechs
Charaktere $\chifn_8,\chifn_{13},\chifn_{21},\chifn_{29},\chifn_{33},
\chifn_{60}$); dies ist \emph{keine} Verifikation der Formel, sondern
eine in \S\ref{sec:tautology} aufgedeckte Tautologie.

Zweitens erreicht der Prädiktor $S^{(600)}+\archdiff^{(600)}$ eine hohe
lineare Korrelation von $R^2=0.95$ gegen die empirischen Gaps, fällt aber auf
eine Vorzeichengenauigkeit von $8/10$ aufgrund von zwei Vorzeichenwechseln bei $\chifn_{21}$
und $\chifn_{29}$.

\subsection{Regressionsstatistik}
\label{sec:validation-regression}

Wir halten alle vier kanonischen Regressionskonfigurationen in
Tabelle~\ref{tab:regression} fest. Das Muster ist für das Folgende von zentraler Bedeutung.

\begin{table}[t]
\caption{Vorzeichen- und lineare Regressions-Performance der
  Prädiktor-Realisierungen gegen den Galerkin-Referenz-Gap
  $\gap_{\mathrm{gal}}^{(\mathrm{ref})}$ bei $\lambda=20000$,
  $N\in\{200,600\}$. Die Vorzeichengenauigkeit wird relativ zum
  Vorzeichen des hoch-$N$ Galerkin-Referenz-Gaps in der
  Connes-Konvention berechnet (Bemerkung~\ref{rem:sign-convention-dependence}).
  $^*$ Die $N=200$-Statistiken unterliegen den
  Galerkin-Trunkierungseffekten, die in
  Bemerkung~\ref{rem:sw-compression} und
  \S\ref{sec:discussion-slope} dokumentiert sind; die $N=600$-Zeile ist
  das operationell vorzuziehende Resultat.
  Der Ausschluss von $\chifn_{21}$ (dem einzigen Vorzeichenausreißer;
  siehe \S\ref{sec:chi21}) ergibt per Konstruktion eine 9/9
  Vorzeichengenauigkeit bei $N=600$ und wird daher nicht als
  eigenständige Statistik berichtet.}
\label{tab:regression}
\centering
\small
\begin{tabular}{@{}lccc@{}}
\toprule
Konfiguration & $N$ & sign\_ok & $R^2$ \\
\midrule
$\gap_{\mathrm{gal}}$ (alle 10)$^*$ & $200$ & $9/10$ & $0.41$ \\
$S+\archdiff$ (alle 10)$^*$ & $200$ & $9/10$ & $0.42$ \\
\midrule
$\gap_{\mathrm{gal}}$ (alle 10) & $600$ & $10/10$ & $0.99$ \\
$S+\archdiff$ (alle 10) & $600$ & $8/10$ & $0.95$ \\
\bottomrule
\end{tabular}

\smallskip
\footnotesize
Die Vorzeichengenauigkeit ist relativ zur durchgängig verwendeten
Connes-Operator-Konvention $\tilde{A}_\chi$ berechnet
(Anhang~\ref{app:sign}). Die entgegengesetzte Konvention
(Weil-$Q$, Koeffizient $+2$) würde alle Vorzeichen invertieren und
eine komplementäre Vorzeichenstatistik liefern, jedoch bleibt die
relative Übereinstimmung Prädiktor-vs.-Referenz unverändert
(vgl.\ Bemerkung~\ref{rem:sign-convention-dependence}).
\end{table}

Bei $N=200$ erreichte der Prädiktor $S+\archdiff$ nach Ausschluss von
$\chifn_{21}$ eine scheinbare $9/9$ Vorzeichengenauigkeit mit $R^2=0.80$
--- die numerische Schließung, die den Titel des „Weil-kernel closure
paper" in der frühen Entwicklungs-Session 7 motivierte. Bei $N=600$
passieren jedoch drei Dinge gleichzeitig:
(i) Der Galerkin-Gap $\gap_{\mathrm{gal}}^{(600)}$ wird
numerisch ununterscheidbar von der empirischen Referenz für sechs
der zehn Charaktere (Tautologie, \S\ref{sec:tautology});
(ii) Der Galerkin-Gap von $\chifn_{21}$ kollabiert von $+0.282$ auf
$-0.004$, was ihn als einen $N$-Oszillator bestätigt
(\S\ref{sec:chi21}); und
(iii) Der Prädiktor $S+\archdiff^{(600)}$ gewinnt einen Punkt bei
$R^2$, verliert aber zwei Punkte bei der Vorzeichengenauigkeit und scheitert nun
sowohl bei $\chifn_{21}$ (wo der $N=600$ Galerkin-Gap klein ist) als auch bei
$\chifn_{29}$ (wo Realisierung (R1) im Vorzeichen von
Realisierung (R2) abweicht).

\begin{remark}[Siegel-Walfisz-Kompression bei $N=600$]
\label{rem:sw-compression}
Der Vergleich der $\gap_{\mathrm{gal}}$-Spalten bei $N=200$ und $N=600$
offenbart einen systematischen Kompressionseffekt: Für den stark-positiven
Charakter $\chifn_{12}$ fällt $\gap_{\mathrm{gal}}$ von
$+0.0325$ auf $+0.0115$ (Faktor $0.35$); für $\chifn_{17}$,
$\chifn_{24}$, $\chifn_{29}$, $\chifn_{60}$ gilt derselbe
Schrumpfungsfaktor bis auf $\pm 0.1$. Die stark-negativen
Charaktere $\chifn_8$, $\chifn_{13}$, $\chifn_{33}$ sind
qualitativ stabil zwischen $N=200$ und $N=600$ (relative Änderung
$\lesssim 10\%$). Diese Dichotomie erinnert qualitativ an die obere
Siegel-Walfisz-Schranke
$|\sum_{p\leq x}\chifn(p)\log p|=o(\sqrt{x})$, angewandt auf den
Schwanz der großen Primzahlen --- Charaktere mit einem echten
Signal bei kleinem Führer erscheinen weitgehend bereits bei $N=200$ bestimmt,
während schwach signalisierende Charaktere ihre bei $N=200$ ermittelten
Schätzungen bei $N=600$ um einen Faktor von etwa $0.35$ komprimiert sehen
--- aber es wird hier keine theoretische Herleitung angeboten, die den
Kompressionsfaktor mit einer Siegel-Walfisz-Konstante verknüpft. Wir
markieren die empirische Beobachtung als offene Frage; siehe
\S\ref{sec:outlook}, Punkt zur theoretischen Herleitung der
Siegel-Walfisz-Kompression.
\end{remark}

\begin{table*}[t]
\caption{$N$-Konvergenz des Galerkin-Gaps und des Weil-Kern-Prädiktors
  bei $\lambda=20000$ für zehn reelle Charaktere mit kleinem Führer. Beachte, dass
  die $\gap_{\mathrm{gal}}^{(\mathrm{ref})}$-Einträge bei $N=600$ identisch mit
  $\gap_{\mathrm{gal}}^{(N=600)}$ übereinstimmen, da die empirischen Referenzwerte
  selbst Ausgaben derselben Galerkin-Diagonalisierung bei $N=600$ sind;
  dies ist die Tautologie, die in \S\ref{sec:tautology} diskutiert wird.
  Quellen: \texttt{ARCH\_TERM\_ANALYSIS.md}, \texttt{ARCH\_TERM\_N600\_ANALYSIS.md}.}
\label{tab:validation}
\centering
\small
\begin{tabular}{@{}lrrrrrr@{}}
\toprule
$\chifn$ & $D$ & $\gap_{\mathrm{gal}}^{(\mathrm{ref})}$ ($N_\mathrm{src}$) & $\gap_{\mathrm{gal}}^{(200)}$ & $\gap_{\mathrm{gal}}^{(600)}$ & $(S+\archdiff)^{(200)}$ & $(S+\archdiff)^{(600)}$ \\
\midrule
$\chi_5$  &  5 & $+0.01560\ (400)$ & $+0.01429$ & $+0.01603$ & $+3.35\!\cdot\!10^{-2}$ & $+3.37\!\cdot\!10^{-2}$ \\
$\chi_8$  &  8 & $-0.04902\ (600)$ & $-0.04548$ & $-0.04902$ & $-4.22\!\cdot\!10^{-2}$ & $-4.08\!\cdot\!10^{-2}$ \\
$\chi_{12}$ & 12 & $+0.03367\ (400)$ & $+0.03251$ & $+0.01154$ & $+5.07\!\cdot\!10^{-2}$ & $+2.03\!\cdot\!10^{-2}$ \\
$\chi_{13}$ & 13 & $-0.12504\ (600)$ & $-0.12502$ & $-0.12504$ & $-1.15\!\cdot\!10^{-1}$ & $-1.16\!\cdot\!10^{-1}$ \\
$\chi_{17}$ & 17 & $+0.01318\ (400)$ & $+0.00408$ & $+0.00886$ & $+1.67\!\cdot\!10^{-2}$ & $+1.13\!\cdot\!10^{-2}$ \\
$\chi_{21}$ & 21 & $-0.00423\ (600)$ & $\mathbf{+0.28163}$ & $-0.00423$ & $+2.89\!\cdot\!10^{-1}$ & $+1.14\!\cdot\!10^{-2}$\,\textbf{(F)} \\
$\chi_{24}$ & 24 & $+0.01076\ (400)$ & $+0.01632$ & $+0.00866$ & $+1.74\!\cdot\!10^{-2}$ & $+1.86\!\cdot\!10^{-2}$ \\
$\chi_{29}$ & 29 & $+0.01439\ (600)$ & $+0.02372$ & $+0.01439$ & $+2.58\!\cdot\!10^{-2}$ & $-1.16\!\cdot\!10^{-2}$\,\textbf{(F)} \\
$\chi_{33}$ & 33 & $-0.14221\ (600)$ & $-0.14051$ & $-0.14221$ & $-1.31\!\cdot\!10^{-1}$ & $-1.32\!\cdot\!10^{-1}$ \\
$\chi_{60}$ & 60 & $+0.00521\ (600)$ & $+0.12363$ & $+0.00521$ & $+1.27\!\cdot\!10^{-1}$ & $+7.21\!\cdot\!10^{-3}$ \\
\midrule
\multicolumn{4}{l}{\textit{$S+\archdiff$ Vorzeichengenauigkeit ($N=200$):}} & & $9/10$ ($R^2=0.42$) & \\
\multicolumn{4}{l}{\textit{$S+\archdiff$ Vorzeichengenauigkeit ($N=600$):}} & & & $8/10$ ($R^2=0.95$) \\
\multicolumn{4}{l}{\textit{$\gap_{\mathrm{gal}}$ Vorzeichengenauigkeit ($N=600$):}} & & & $10/10$ ($R^2=0.99$, Tautologie) \\
\bottomrule
\end{tabular}
\end{table*}

\subsection{Evidenzschichten der Tabellen}
\label{sec:validation-evidence-ledger}

Die numerische Präsentation ist als geschichtete Diagnose zu lesen, nicht
als eine einzelne Validierungsstatistik. Tabelle~\ref{tab:evidence-layer-ledger}
macht explizit, welche Schlussfolgerung jede Schicht tragen kann und
welche Schlussfolgerung außerhalb der vorliegenden Evidenzbasis bleibt.

\begin{table*}[t]
\caption{Evidenzschichten-Ledger für den numerischen Atlas. Die Tabelle
  hält fest, was jede Datenschicht legitimerweise trägt und welche
  Schlussfolgerung außerhalb der vorliegenden Evidenzbasis bleibt.}
\label{tab:evidence-layer-ledger}
\centering
\footnotesize
\begin{tabular}{@{}p{0.16\textwidth}p{0.22\textwidth}p{0.28\textwidth}p{0.24\textwidth}@{}}
\toprule
Schicht & Datenanker & Trägt & Trägt nicht \\
\midrule
Roh-Gap-Schicht &
Tabelle~\ref{tab:validation} und die dort genannten Quellberichte &
Hoch-$N$-Galerkin-Selbstkonsistenz und die Lokalisierung
vorzeichensensitiver Charaktere wie $\chifn_{21}$, $\chifn_{29}$ und
$\chifn_{60}$. &
Externe Wahrheit des $N\to\infty$-Gaps oder Schließung von
$C^{(2)}_{\chifn}$. \\
\addlinespace
Regressionsschicht &
Tabelle~\ref{tab:regression} &
Die scheinbare Schließung bei $N=200$ ist trunkierungssensitiv; der
$N=600$-Prädiktor hat hohes $R^2$, aber zwei Vorzeichenfehler. &
Eine universelle Vorzeichenregel oder eine konventionsfreie
Validierungsstatistik. \\
\addlinespace
Atlas-Taxonomieschicht &
Tabelle~\ref{tab:archratio} und
\S\ref{sec:gapsig} &
Stichprobenspezifische Prim-/Mixed-/Arch-Regime und
Gap-Signaturklassen. &
Universalität über die zehn Charaktere mit kleinem Führer hinaus. \\
\addlinespace
Failure-Control-Schicht &
Proposition~\ref{prop:tautology} und
Tabelle~\ref{tab:chi21} &
Die Matrixzerlegungsidentität und der $\chifn_{21}$-Oszillator blockieren
eine Schließungsdeutung der Galerkin-Numerik. &
Eine ersetzende analytische Theorie der Grundzustände. \\
\addlinespace
Fehlende-Referenz-Schicht &
\S\ref{sec:limitation-external} &
Benennung der externen Gates vor prädiktiven Claims: höherpräzise
Galerkin-Rechnung, LMFDB-Nullstellenvergleich oder unabhängige Basis. &
Eine bereits abgeschlossene externe Validierung im vorliegenden Atlas. \\
\bottomrule
\end{tabular}
\end{table*}
\section{Charakter-Atlas in vier Achsen}
\label{sec:atlas}

Die Rohtabelle aus Abschnitt~\ref{sec:validation} erlaubt eine strukturelle
Lesart entlang vier taxonomischer Dimensionen: zwei \emph{Filter-Konstanten}
(Parität und Wurzelzahl, beide gleich $+1$ im gesamten Sample) und zwei
\emph{informative Achsen} (archimedisches Gewichtsverhältnis und Gap-Signatur).
Wir entfalten jede Dimension nacheinander, in der Reihenfolge zunehmender
Aussagekraft.

\subsection{Paritäts-Kartografie}
\label{sec:parity}

Die Parität $\chifn(-1)\in\{+1,-1\}$ entscheidet über den archimedischen Faktor von
$\Lambda(s,\chifn)$: Gerade Charaktere verwenden $\Gamma(s/2)$, ungerade
Charaktere verwenden $\Gamma((s+1)/2)$. In dieser Arbeit sind alle zehn Charaktere
Kronecker-Charaktere $\chifn_D$ mit positiver Diskriminante $D$, und
daher alle gerade. Dies ist Absicht: Die durchgehend verwendete Sektor-Zerlegung
$H_L=H_L^+\oplus H_L^-$ setzt eine gerade Parität voraus, welche
eine nicht-triviale Kosinus-Sinus-Asymmetrie des anti-diagonalen Kerns
nach Theorem~\ref{thm:sector-diff} sicherstellt (ungerade Charaktere würden
den Paired-Space-Formalismus (paired-space formalism)
$\chifn\oplus\bar\chifn$ oder einen Tensor mit einem Vorzeichen-Twist erfordern; vgl.\
\cite{MetaGalerkinBridge}, \S5.5, Fußnote zur Gamma-Parität). Die
Beschränkung auf gerade Charaktere fixiert somit das qualitative Regime und lässt Raum
für die drei nachfolgenden, feineren Klassifikationen.

\subsection{Wurzelzahl-Klassifikation}
\label{sec:rootnumber}

Für jeden reellen Dirichlet-Charakter ist die Wurzelzahl $W(\chifn):=
\epsilon_{\chifn}$ durch die klassische Gaußsche Summenformel
$\tau(\chifn)=\sqrt{q}$ für primitive reelle gerade Charaktere \cite{IwaniecKowalski}
gleich $+1$. Alle zehn Charaktere in unserem Atlas teilen somit
dieselbe Wurzelzahl $W=+1$, und die Wurzelzahl ist über die gesamte Tabelle
hinweg konstant. Folglich kann die Wurzelzahl Gap-Vorzeichen
innerhalb unserer Stichprobe nicht diskriminieren. Eine größere Stichprobe, die sich auf komplexe
Charaktere erstreckt (\S\ref{sec:outlook}), würde eine echte Wurzelzahl-Achse
mit $W\in\{+1,-1,\pm i\}$ aufzeigen; wir markieren dies als Teil des offenen
Ausblicks und nicht als ein Ergebnis des vorliegenden Atlas.

\subsection{Schicht des archimedischen Faktors}
\label{sec:archlayer}

Die Größe
$r_{\chifn}:=|\archdiff_{\chifn}^{(200)}|/|S_{\chifn}^{(200)}|$
(Verhältnis von archimedischem Beitrag zu Primzahlbeitrag bei $N=200$,
$\lambda=20000$) variiert über die Stichprobe hinweg um fast zwei Größenordnungen.
Tabelle~\ref{tab:archratio} dokumentiert diese Werte.

\begin{table}[t]
\caption{Archimedisches Gewichtsverhältnis (Archimedean weight ratio)
  $r_{\chifn}=|\archdiff|/|S|$ bei $N=200$, $\lambda=20000$.}
\label{tab:archratio}
\centering\scriptsize
\begin{tabular}{@{}lcccc@{}}
\toprule
$\chifn_D$ & $|S_{\chifn}|$ & $|\archdiff_{\chifn}|$ & $r_{\chifn}$ & Klasse \\
\midrule
$\chifn_{29}$ & $2.54\!\cdot\!10^{-2}$ & $4.47\!\cdot\!10^{-4}$ & $1.8\%$ & prim-dominiert \\
$\chifn_{24}$ & $1.62\!\cdot\!10^{-2}$ & $1.19\!\cdot\!10^{-3}$ & $7.4\%$ & prim-dominiert \\
$\chifn_{33}$ & $1.65\!\cdot\!10^{-1}$ & $3.44\!\cdot\!10^{-2}$ & $21\%$ & prim-dominiert \\
$\chifn_5$ & $9.06\!\cdot\!10^{-2}$ & $5.71\!\cdot\!10^{-2}$ & $63\%$ & gemischt \\
$\chifn_8$ & $1.16\!\cdot\!10^{-1}$ & $7.43\!\cdot\!10^{-2}$ & $64\%$ & gemischt \\
$\chifn_{13}$ & $1.02\!\cdot\!10^{-1}$ & $1.33\!\cdot\!10^{-2}$ & $13\%$ & prim-dominiert \\
$\chifn_{17}$ & $2.02\!\cdot\!10^{-2}$ & $3.52\!\cdot\!10^{-3}$ & $17\%$ & prim-dominiert \\
$\chifn_{12}$ & $6.51\!\cdot\!10^{-2}$ & $1.44\!\cdot\!10^{-2}$ & $22\%$ & prim-dominiert \\
$\chifn_{21}$ & $2.31\!\cdot\!10^{-1}$ & $5.77\!\cdot\!10^{-2}$ & $25\%$ & prim-dominiert$^\dagger$ \\
$\chifn_{60}$ & $4.95\!\cdot\!10^{-2}$ & $1.76\!\cdot\!10^{-1}$ & $356\%$ & arch-dominiert \\
\bottomrule
\end{tabular}

\smallskip
\footnotesize
$^\dagger$ bei $N=200$; die effektiven Gewichte verschieben sich bei $N=600$,
während sich die Galerkin-Grundzustände stabilisieren (vgl.\ \S\ref{sec:chi21}).
\end{table}

Drei Unterklassen kristallisieren sich heraus:
\begin{itemize}
  \item \emph{Prim-dominiert} ($r_{\chifn}\leq 25\%$):
    $\chifn_{12}$, $\chifn_{13}$, $\chifn_{17}$, $\chifn_{21}$,
    $\chifn_{24}$, $\chifn_{29}$, $\chifn_{33}$. Die Primsumme
    $S_{\chifn}$ bestimmt sowohl die Größe als auch das Vorzeichen des Gaps;
    $\archdiff_{\chifn}$ verschiebt den quantitativen Wert um nicht mehr
    als ein Viertel. Dies ist das „natürliche“ Regime, in dem die
    Intuition bezüglich des Weil-Kerns greift.
  \item \emph{Gemischt} ($r_{\chifn}\in[60\%,70\%]$):
    $\chifn_5$, $\chifn_8$. Primzahlen und archimedischer Faktor sind
    vergleichbar; keines von beidem ist isoliert betrachtet ein verlässlicher Prädiktor.
  \item \emph{Arch-dominiert} ($r_{\chifn}>100\%$):
    $\chifn_{60}$. Hier \emph{übersteigt} die archimedische Sektordifferenz
    den Primbeitrag um einen Faktor $\sim 4$, und
    $S_{\chifn}$ sowie $\archdiff_{\chifn}$ unterscheiden sich sogar im Vorzeichen
    ($S=-4.95\!\cdot\!10^{-2}$, $\archdiff=+1.76\!\cdot\!10^{-1}$).
    Der empirische Gap $\gap_{\mathrm{gal}}^{(\mathrm{ref})}=+0.005$ ist nahezu null und
    wird vollständig vom archimedischen Term getragen.
\end{itemize}

Die Unterklassen-Unterscheidung ist struktureller, nicht nur quantitativer Natur: Für
prim-dominierte Charaktere würde sich eine künftige analytische Kontrolle von
$\Deltachi(m\log p)$ direkt in eine Vorzeichen-Vorhersage des Gaps übersetzen,
wohrend für arch-dominierte Charaktere die archimedische
$\psi$-Funktionsasymptotik dominiert und eine unabhängige Analyse erforderlich ist.

\begin{remark}[Deskriptiver vs.\ prädiktiver Status der Trichotomie]
\label{rem:trichotomy-descriptive}
Mit zehn Charakteren ist die Prim/Mixed/Arch-Trichotomie von
Tabelle~\ref{tab:archratio} deskriptiv für die Stichprobe, nicht
prädiktiv für beliebige primitive reelle Charaktere: Die Schwellenwerte
($r\leq 25\%$, $r\in[60\%,70\%]$, $r>100\%$) wurden so gewählt, dass
sie die beobachteten Cluster trennen, mit keinem Datenpunkt in den
Lücken-Intervallen $r\in(25\%,60\%)$ und $r\in(70\%,100\%)$. Ob die
Trichotomie verallgemeinert, ist eine Hypothese, die an der
Leiter-Erweiterung $D\leq 100$ (\S\ref{sec:outlook}, Punkt~4) zu testen
wäre; der vorliegende Atlas behauptet sie nicht als universelle
Klassifikation.
\end{remark}

\subsection{Gap-Signatur-Atlas}
\label{sec:gapsig}

Die folgenden vier Klassen ordnen die zehn Charaktere anhand der
quantitativen Natur ihrer Gaps, wobei $\gap_{\mathrm{gal}}^{(\mathrm{ref})}$ aus
Tabelle~\ref{tab:validation} abgelesen wird.

\paragraph{Stark-negativer Prototyp.}
$\chifn_{33}$ ist der herausragende Fall:
$\gap_{\mathrm{gal}}^{(\mathrm{ref})}=-0.142$, mit nahezu identischen Werten bei
$N\in\{200,400,600\}$. Sowohl der Primbeitrag als auch der empirische Gap
sind durchweg negativ, und der Gap ist absolut betrachtet der größte
in der Stichprobe. Der Charakter ist
$\chifn_{33}(n)=\bigl(\tfrac{33}{n}\bigr)$, mit der Faktorisierung
$33=3\cdot 11$; die entsprechenden Legendre-Symbol-Werte bei kleinen
Primzahlen sind $\chifn_{33}(2,5,7,13,17,19,23,29,31,37)=
(1,-1,-1,-1,-1,-1,+1,-1,+1,-1)$, was einen klaren Bias (bias) in Richtung
$-1$ liefert. Wir bezeichnen $\chifn_{33}$ als den \emph{stark-negativen
Prototyp} des Atlas.

\paragraph{Stark-positiver Prototyp.}
$\chifn_{12}$ ist das Pendant auf der positiven Seite:
$\gap_{\mathrm{gal}}^{(\mathrm{ref})}=+0.034$ bei $N=400$, komprimiert auf $+0.012$ bei
$N=600$ (Bemerkung~\ref{rem:sw-compression}). Der primitive Charakter
mod $12$ besitzt als Führer ein Produkt aus zwei kleinen Primzahlen ($q=12=
4\cdot 3$); frühere asymptotische Studien \cite{AsymptoticScan} beobachteten
über drei Größenordnungen in $\lambda$ hinweg bei $N=200$ eine Steigung (slope) konsistent mit null,
was wir nun als eine durch $N$ limitierte
Beobachtung re-interpretieren (siehe auch \texttt{THEORIE\_CHI12\_KONSTANZ.md}). Innerhalb
des vorliegenden Atlas ist $\chifn_{12}$ der stark-positive Prototyp,
wenngleich mit geringerer Ausprägung als $\chifn_{33}$.

\paragraph{Moderat-negative Klasse.}
$\chifn_8$ und $\chifn_{13}$ zeigen beide stabile negative Gaps
($-0.049$ bzw.\ $-0.125$) mit konsistenten Vorzeichen über
$N\in\{200,400,600\}$ hinweg. Sie liegen zwischen dem stark-negativen
Prototyp und dem schwach-positiven Cluster.

\paragraph{Schwach-positiver Cluster.}
$\chifn_{17}$, $\chifn_{24}$, $\chifn_{29}$ bilden einen engen Cluster
kleiner positiver Gaps ($+0.013$, $+0.011$, $+0.014$). Ihre absoluten
Werte liegen innerhalb des Siegel-Walfisz-Kompressionsregimes; ihr
Prädiktor $S+\archdiff$ behält bei $N=200$ sein Vorzeichen, kann jedoch bei
$N=600$ kippen (Fall von $\chifn_{29}$, Realisierung (R1); vgl.\
\S\ref{sec:validation-regression}). Dieser Cluster markiert die untere
Schranke der zuverlässigen numerischen Vorzeichenauflösung für unsere aktuellen
Methoden.

\paragraph{Arch-dominierter Fast-null-Fall.}
$\chifn_{60}$ hat $\gap_{\mathrm{gal}}^{(\mathrm{ref})}=+0.005$ mit
$r_{\chifn_{60}}=356\%$ (s.o.). Der Prädiktor $S+\archdiff$ besitzt bei $N=600$ zufällig
das korrekte Vorzeichen, aber nicht aufgrund
des Primbeitrags, welcher im Vorzeichen dem empirischen
Gap widerspricht; stattdessen dominiert der archimedische Term. Wir kennzeichnen $\chifn_{60}$
als ein distinktes Regime, das eine eigene Analyse rechtfertigt.

\paragraph{Composite-Conductor-Oszillator.}
$\chifn_{21}$ ($D=21=3\cdot 7$) hat $\gap_{\mathrm{gal}}^{(\mathrm{ref})}=-0.004$ bei
$N=600$, aber einen $N=200$ Galerkin-Gap von $+0.282$ --- ein Vorzeichenwechsel
um den Faktor $66$. Er ist der einzige Charakter in unserer Stichprobe, der dieses
Verhalten zeigt, und wir behandeln ihn detailliert in \S\ref{sec:chi21}.

\smallskip

Die obige Vier-Achsen-Kartografie ist der deskriptive Ertrag dieser Arbeit:
Sie gruppiert die zehn Charaktere in qualitativ unterschiedliche Regime, und
zwar bevor und unabhängig von jeder prädiktiven Formel. Insbesondere
zeigen die Trichotomie prim-dominiert/gemischt/arch-dominiert
(Tabelle~\ref{tab:archratio}) und die Gap-Signatur-Klassen
(stark-neg, stark-pos, moderat-neg, schwach-pos, arch-dominiert,
Oszillator) zusammengenommen sechs Unterklassen bei zehn Charakteren auf --- ein
erstes Anzeichen dafür, dass der „Ein-Achsen-Transfer“ (single-axis transfer) des klassischen
Regimes der geraden Dominanz \cite{RH_II} zu grob ist.

\section{Die Matrixzerlegungs-Tautologie}
\label{sec:tautology}

\begin{proposition}[Tautologie]
\label{prop:tautology}
Seien $\phi^\pm_{\chifn}$ die Grundzustands-Eigenvektoren der Galerkin-trunkierten
Operatoren $K^\pm_{\chifn}$ bei Basisgröße $N$. Dann gilt für jede Zerlegung
$W^\pm = W^\pm_{\mathrm{arch}} + W^\pm_{\mathrm{prime}}$ der Galerkin-Matrix:
\[
\begin{aligned}
  \gap_{\mathrm{gal}}^{(N)}
    &= \archdiff + \primediff,\\
  \archdiff
    &:= \langle\phi^-, W^-_{\mathrm{arch}}\phi^-\rangle
      - \langle\phi^+, W^+_{\mathrm{arch}}\phi^+\rangle,\\
  \primediff
    &:= \langle\phi^-, W^-_{\mathrm{prime}}\phi^-\rangle
      - \langle\phi^+, W^+_{\mathrm{prime}}\phi^+\rangle.
\end{aligned}
\]
\end{proposition}

\begin{proof}
Die Linearität des inneren Produkts auf dem Galerkin-Raum ergibt
$\langle\phi^\pm, W^\pm\phi^\pm\rangle = \langle\phi^\pm,
W^\pm_{\mathrm{arch}}\phi^\pm\rangle + \langle\phi^\pm,
W^\pm_{\mathrm{prime}}\phi^\pm\rangle$.
Da $\phi^\pm$ Eigenvektoren mit Eigenwert $\lambda_1^{\min}(W^\pm)$ sind,
zerfällt der Gap $\gap_{\mathrm{gal}}^{(N)}$ als Differenz der beiden rechten
Seiten. $\square$
\end{proof}

Obwohl der propositionale Gehalt eine triviale Linearitätsidentität ist, ist
die \emph{methodologische Konsequenz} nicht-trivial: Jede scheinbare
Übereinstimmung zwischen $\gap_{\mathrm{gal}}^{(N)}$ und $S_\chi+\archdiff_\chi$
bei festem $N$ ist strukturell, keine Vorhersage. Wenn
$\mathrm{prime\_diff}_{\mathrm{matrix}}$ durch die Autokorrelationsformel
$S_{\chifn}$ in der Raumdomäne ersetzt wird, entsteht eine kleine numerische Diskrepanz
(in der Größenordnung $10^{-3}$, etwa $1$--$2\%$ von $S$); die „9/9 Vorzeichengenauigkeit bei
$N=200$", von der in Vorarbeiten des Autors berichtet wurde, lässt zwei
gleichermaßen konsistente Lesarten zu: (a) Die Diskrepanz war für die
Stichprobe von zehn Charakteren bei niedrigem $N$ \emph{glücklich klein}
(fortuitous smallness) und wächst bei höherem $N$, sobald die intrinsischen
Limitationen des Prädiktors aufgelöst werden; (b) Bei $N=600$ setzt für
Charaktere mit $|\gap|\lesssim 10^{-2}$ eine
\emph{Konditionierungspathologie} ein, bei der Gleitkomma- und
Interpolationsfehler der Größenordnung $10^{-3}$ vorzeichenrelevant
werden. Die zwei Lesarten sind ohne externe Referenz
(vgl.\ \S\ref{sec:limitation-external}) nicht zu unterscheiden. In
beiden Fällen ist die $N=200$-Übereinstimmung keine strukturelle
Eigenschaft der Formel. Bei $N=600$ wird die Diskrepanz groß genug,
um Vorzeichen für $\chifn_{29}$ umzukehren (siehe Tabelle~\ref{tab:validation}).

\begin{remark}[Warum das autokorrelations-numerische $S$ scheitert]
Die Autokorrelation $(\phi_{\chifn}^+\ast\phi_{\chifn}^+)(t)$ erfordert
die numerische Auswertung auf einem diskreten Raster (grid), kombiniert mit einer Interpolation auf
Primzahlabzissen $t=m\log p$. Gleitkomma- und Interpolationsfehler
auf der Rasterskala $\Delta x = 2L/n_{\mathrm{grid}}$ übersetzen sich in Fehler
der Größenordnung $10^{-3}$ in $S$, was für schwach signalisierende
Charaktere wie $\chi_{21},\chi_{29}$ die gleiche Größenordnung wie $|\gap|$ darstellt.
\end{remark}

\begin{corollary}[Woher die prädiktive Kraft kommen muss]
\label{cor:paleywiener}
Eine wahrhaft prädiktive Weil-Kern-Gap-Formel erfordert die analytische Kontrolle
der Grundzustände $\phi^\pm_{\chifn}$ \emph{ohne} den Weg über
eine Galerkin-Diagonalisierung. Die Route ist eine asymptotische Paley-Wiener-Analyse
(\`a la Connes \cite{Connes}, angepasst auf das Dirichlet-Setting) oder eine
direkte spektrale Identifikation von $\Deltachi(t)$ aus den lokalen Daten
von $\chifn$ (Parität, Wurzelzahl, Führer, arithmetischer Faktor der
$L$-Funktion).
\end{corollary}

\section{\texorpdfstring{Die $\chi_{21}$ $N$-Oszillation}{Die chi21 N-Oszillation}}
\label{sec:chi21}

Der Charakter $\chifn_{21}$, mit zusammengesetztem Führer $q=21=3\cdot 7$,
ist der einzige Eintrag aus Tabelle~\ref{tab:validation}, bei dem der Galerkin-Gap
bei $N=200$ und die empirische Referenz bei $N=600$ im Vorzeichen abweichen.
Wir quantifizieren die Diskrepanz und führen sie auf einen
Galerkin-Trunkierungs-Effekt zurück, der bei hohem $N$ aufgelöst wird, nicht auf einen
Defekt der Formel.

Tabelle~\ref{tab:chi21} dokumentiert die $N$-Konvergenzdaten bei
$\lambda=20000$, aus dem Skript \texttt{chi21\_n\_convergence.py}.

\begin{table}[h]
\caption{$N$-Konvergenz von $\chifn_{21}$ bei $\lambda=20000$.
  Der Galerkin-Gap ist zwischen $N=200$ und $N=400$ flach und
  springt dann um mehr als eine Größenordnung zwischen $N=400$ und
  $N=600$. Zum Vergleich konvergiert der stark-positive Charakter
  $\chifn_{12}$ monoton (letzte Spalte).}
\label{tab:chi21}
\centering\small
\begin{tabular}{@{}cccc@{}}
\toprule
$N$ & $\gap_{\mathrm{gal}}(\chifn_{21})$ & $S+\archdiff$ & $\gap_{\mathrm{gal}}(\chifn_{12})$ \\
\midrule
$200$ & $+0.2816$ & $+0.2888$ & $+0.0325$ \\
$400$ & $+0.2791$ & $+0.2869$ & $+0.0337$ \\
$600$ & $-0.0042$ & $+0.0114$ & $+0.0115$ \\
\bottomrule
\end{tabular}
\end{table}

Zwei Beobachtungen strukturieren die Analyse.

\begin{itemize}
  \item[(i)] \emph{Stufendiskontinuität (Step discontinuity).} Zwischen $N=200$ und $N=400$
    ist der Galerkin-Gap numerisch stationär
    ($0.2816\to 0.2791$, Änderung $0.9\%$); er fällt dann um einen Faktor
    von $\sim 66$ im absoluten Betrag ab und wechselt das Vorzeichen zwischen $N=400$
    und $N=600$. Dies ist nicht die Signatur einer glatten
    $N^{-\alpha}$ Konvergenz; die Daten sind konsistent mit --- aber
    beweisen nicht eindeutig --- einen
    \emph{vermuteten Quasi-Entartungs-Kollaps (vollständige Konvergenz nur bis $N=600$ verifiziert)},
    bei dem die Galerkin-trunkierten Sektor-Eigenwerte
    $\lambda_1^{\min}(W^\pm_{\chifn_{21},N})$ für $N\in\{200,400\}$
    einem Pseudo-Grundzustand entsprechen, der vom Grundzustand bei
    unendlichem $N$ verschieden ist; der wahre Grundzustand wird erst
    darstellbar, wenn die Basis groß genug ist. Eine definitive
    Diagnose würde spektrale Daten $\lambda_2-\lambda_1$ und
    Eigenvektor-Überlapp $|\langle\phi^{(N=200)},\phi^{(N=600)}\rangle|$
    bei intermediären $N\in\{500,700,800\}$ erfordern; wir markieren
    dies als die direkteste weiterführende Verifikation
    (vgl.\ \S\ref{sec:outlook}, Punkt zur spektralen Diagnose von
    $\chifn_{21}$).
  \item[(ii)] \emph{Der Prädiktor vollzieht den Kollaps nach.} Der Prädiktor
    $S+\archdiff$ wird von der Grundzustands-Autokorrelation dominiert
    und kollabiert daher \emph{ebenfalls} zwischen $N=400$ und $N=600$,
    von $+0.287$ auf $+0.011$. Bei $N=600$ stimmt der Prädiktor im Vorzeichen nicht mehr
    mit dem empirischen Gap überein, aber seine Größenordnung ist nun
    bis auf einen Faktor $\sim 3$ mit der empirischen Größenordnung konsistent.
    Dies bestätigt, dass die Formel selbst nicht
    anomal ist: Sobald der Galerkin-Grundzustand konvergiert ist,
    konvergiert der Prädiktor mit ihm, wobei das Vorzeichen durch die
    numerischen Residuen nahe null bestimmt wird.
\end{itemize}

Als Kontrast zeigt $\chifn_{12}$ (Tabelle~\ref{tab:chi21}, letzte Spalte)
eine monotone $N$-Konvergenz ohne Vorzeichenwechsel:
$+0.0325\to+0.0337\to+0.0115$. Dies ist das generische Verhalten bei den
stabilen Charakteren. Die Diskontinuität für $\chifn_{21}$ ist folglich
weder ein Galerkin-Artefakt, das Composite-Conductor-Charakteren im Allgemeinen innewohnt
(da $\chifn_{12}$, mit $q=12=4\cdot 3$,
zusammengesetzt ist), noch eine systematische Grenze der Methode, sondern eine
vermutete charakter-spezifische Beinahe-Entartung (near-degeneracy),
deren definitive spektrale Analyse (über $\lambda_2-\lambda_1$-Scan
und Eigenvektor-Überlapp) den Rahmen des vorliegenden Atlas sprengt;
siehe \S\ref{sec:outlook}.

\begin{remark}[Kein Formel-Defekt]
\label{rem:chi21-not-defect}
Der Vorzeichenwechsel in Tabelle~\ref{tab:chi21} ist eine Eigenschaft des
Galerkin-Operators $W^\pm_{\chifn_{21},N}$ bei niedrigem $N$, nicht von der
Formel \eqref{eq:gap-unreduced}. Da $S+\archdiff$ in
Begriffen der Galerkin-Eigenvektoren definiert ist, wird jeder
falsche Eigenwert bei niedrigem $N$ auf den Prädiktor übertragen. Das richtige theoretische
Objekt ist $\gap_{\mathrm{gal}}^{(\mathrm{ref})}(\chifn_{21})=-0.00423$ bei $N=600$,
und die richtige Frage bezüglich $\chifn_{21}$ ist, ob eine
Paley-Wiener-Analyse des Grundzustands bei unendlichem $N$ diesen Wert direkt reproduzieren würde,
ohne das Plateau bei $N=200,400$ passieren zu müssen.
\end{remark}

\section{Diskussion}
\label{sec:discussion}

\subsection{Implikationen für die Dreikomponenten-Zerlegung}
\label{sec:discussion-three-component}

Die Dreikomponenten-Zerlegung \eqref{eq:three-component} teilt
jedes RH-artige Problem in die Galerkin-Brücke $\mathrm{GBC}$, die
Connes-Stil-Operatoridentifikationskomponente $C^{(2)}$, und den
Hurwitz-Realität-der-Nullstellen-Transfer auf. Für den trivialen Charakter
isoliert \cite{RH_II} die Route über gerade Dominanz zu
$C^{(2)}_{X=\zeta}$ über die Zieltrennung der Sektoren von
$\sqrt{\lambda}$; in der aktuellen Zenodo-Fassung ist dies eine bedingte
Reduktion mit expliziten verbleibenden Voraussetzungen. Für nicht-triviale
reelle $\chifn$ mit kleinem Führer belegt unser Atlas drei strukturelle
Tatsachen über $C^{(2)}_{\chifn}$:

\begin{itemize}
  \item[(I1)] Die Weil-Kern-Architektur (Theoreme~\ref{thm:kernel}
    und~\ref{thm:sector-diff}) ist rigoros und übertragbar. Die
    anti-diagonale Sektordifferenz trägt die volle Charakter-Abhängigkeit,
    wie vorhergesagt.
  \item[(I2)] Bei $\lambda=20000$ (einem einzelnen $\lambda$-Wert) ist der
    Gap $\gap_{\chifn}(\lambda)$ im Absolutbetrag durch $0.15$ über alle
    zehn Charaktere beschränkt, mit charakter-spezifischem Vorzeichen.
    Eine vollständige $\lambda$-Asymptotik ($\Omega(\sqrt{\lambda})$
    vs.\ $O(1)$ vs.\ intermediäres Wachstum) kann aus einer einzigen
    $\lambda$-Messung nicht abgeleitet werden und bleibt zukünftiger
    Arbeit überlassen; vorläufige Resultate für $\chifn_5$ und
    $\chifn_{12}$ über drei Größenordnungen in $\lambda$ bei $N=200$ sind
    konsistent mit $O(1)$-Wachstum (\cite{AsymptoticScan},
    Begleitmaterial), aber ein systematischer $\lambda$-Scan über alle
    zehn Charaktere ist offen. Die grobe
    Hypothese eines „Universaltransfers“, also der Aussage
    „Gerade Dominanz $\Rightarrow$ gerader Sektor gewinnt für alle
    $\chifn$“ ist bei $\lambda=20000$ widerlegt (da sechs von zehn
    Charakteren negative oder fast-null Gaps haben); die
    $\lambda$-asymptotische Widerlegung erfordert den oben genannten
    dedizierten $\lambda$-Scan. (Eine vollständige $\lambda$-Asymptotik
    kann aus einer einzigen Messung bei $\lambda=20000$ nicht abgeleitet
    werden; diese Frage ist offen.)
  \item[(I3)] Die Identifizierung von $C^{(2)}_{\chifn}$ erfordert die analytische
    Kontrolle der Sektor-Grundzustände $\phi^\pm_{\chifn}$
    selbst, nicht bloß ihrer Eigenwerte. Die Galerkin-Diagonalisierung
    liefert die Eigenwerte, produziert aber aufgrund der Tautologie
    von Proposition~\ref{prop:tautology} keine prädiktiven
    Informationen, die nicht ohnehin als Eingang in die Berechnung eingehen.
\end{itemize}

In der Sprache des Meta-Programms von \cite{MetaGalerkinBridge}: Für
nicht-triviale $\chifn$ ist $C^{(2)}_{\chifn}$ eine echte offene
Komponente, und ihr Abschluss ist das zentrale technische Hindernis für einen
GRH-Beweis im Hilbert-P\'olya-Stil im Dirichlet-Fall. Die
analytische Architektur ist hier bereit, einen solchen Abschluss aufzunehmen; was
fehlt, ist die $\chifn$-spezifische asymptotische Analyse von
$\phi^\pm_{\chifn}$.

Nach der Zookeeper-Arbeit \cite{Zookeeper2026} kann diese fehlende
Komponente in der CCM-Sprache noch schärfer gefasst werden. Die Riemann-Zelle
besitzt nun ein unabhängiges Zookeeper-Kernpaket --- Rang-eins-Zerlegung,
Cauchy-Interlacing und Drei-Lemma-Schließung --- neben
dem v2.1-Paket der geraden Dominanz. Der natürliche Dirichlet-Rücktransfer
ist daher Route~D: Konstruiere einen $\chifn$-getwisteten CCM-Operator mit
einem Führer-Rang-eins-Term und einer $\chifn$-gewichteten arithmetischen
Störung, und vergleiche dann dessen Defektnorm mit den Atlas-Kontrollgrößen.
In dieser Lesart sind die Paley-Wiener-Fehlschläge oben keine
Sackgasse; sie identifizieren den ersten konkreten Stresstest für die
$\chifn$-getwistete CCM-Route.

\subsection{Einschränkung: Fehlen einer externen Referenz}
\label{sec:limitation-external}

Alle Gap-Werte in Tabellen~\ref{tab:validation}, \ref{tab:regression}
und~\ref{tab:archratio} entstammen derselben Galerkin-Diagonalisierung
bei $N\in\{200,400,600\}$. Es gibt daher \emph{keine externe,
konventionsfreie Referenz} in diesem Atlas: Die höchsten-$N$-Werte
$\gap_{\mathrm{gal}}^{(\mathrm{ref})}$ sind die bestmöglichen
Galerkin-Schätzungen, keine Messungen des wahren $N\to\infty$-Gaps.
Die statistische Übereinstimmung des Prädiktors $S_\chi+\archdiff_\chi$
mit $\gap_{\mathrm{gal}}^{(\mathrm{ref})}$ ist daher ein
Selbstkonsistenz-Check innerhalb der Galerkin-Methode, keine externe
Validierung. Eine echte externe Referenz würde erfordern:
\begin{itemize}
\item[(i)] mp-Arithmetik-Galerkin bei $N\geq 2000$ zur Diagnose
  endlicher-Präzisions-Artefakte;
\item[(ii)] direkte spektrale Auswertung über LMFDB-Nullstellen von
  $L(s,\chi_D)$;
\item[(iii)] ein unabhängiges Galerkin-Schema mit anderer Basis (z.\,B.\
  Hermite statt Fourier).
\end{itemize}
Keines davon wird hier geliefert; wir markieren dies als offenes
methodologisches Element, bevor prädiktive Behauptungen über
$\gap_{\chifn}(\lambda)$ beim vorliegenden $\lambda$ gestellt werden
können.

\begin{remark}[Konventionsabhängigkeit der Vorzeichenstatistik]
\label{rem:sign-convention-dependence}
Die Vorzeichengenauigkeit in Tabelle~\ref{tab:regression} ist relativ
zur Connes-Operator-Konvention $\tilde{A}_\chi$ berechnet
(Anhang~\ref{app:sign}). Unter der äquivalenten Weil-$Q$-Konvention
invertieren sich alle Gap-Vorzeichen (Prädiktor und Referenz)
gleichzeitig, sodass die relative Übereinstimmung
Prädiktor-vs.-Referenz unverändert bleibt. Die Statistik prüft
folglich Selbstkonsistenz innerhalb einer gewählten Konvention, nicht
ein konventionsfreies Vorzeichen. Eine konventionsfreie Vorzeichen-
Aussage würde eine externe arithmetische Identifikation erfordern
(z.\,B.\ aus der Spektralseite über LMFDB), die wir nicht liefern.
\end{remark}

\subsection{\texorpdfstring{Die Regression mit Steigung $0.72$ und ihre Bedeutung}{Die Regression mit Steigung 0.72 und ihre Bedeutung}}
\label{sec:discussion-slope}

Eine zweite quantitative Tatsache, die durch unsere Daten offengelegt wird, ist die Steigung (slope)
der linearen Regression von $\gap_{\mathrm{gal}}^{(\mathrm{ref})}$ gegen
$S+\archdiff$ bei $N=200$:
$\gap_{\mathrm{gal}}^{(\mathrm{ref})}\approx 0.72\cdot(S+\archdiff)-0.01$
(Stichprobe ohne $\chifn_{21}$, $R^2=0.80$). Wir haben dies in Session~7 der
Entwicklung mit Enthusiasmus festgehalten, doch spätere Reflexion zeigt, dass
sein inhaltlicher Gehalt bescheiden ist. Die Steigung misst die multiplikative
Diskrepanz (mismatch) zwischen dem $N=200$ Prädiktor $S^{(200)}+\archdiff^{(200)}$
und dem empirischen Gap bei hohem $N$. Unter der Siegel-Walfisz-Kompression
aus Bemerkung~\ref{rem:sw-compression} verschwindet diese Diskrepanz
bei höherem $N$: Bei $N=600$ \emph{ist} der Galerkin-Gap
die empirische Referenz (durch Identität, nicht durch Prädiktion), sodass die
Steigung konstruktionsbedingt gegen $1$ geht. Die $0.72$ ist daher eine
$N=200$-spezifische Galerkin-Trunkierungskonstante, keine theoretische
Invariante der Formel.

Ein echtes theoretisches Analogon der Steigung wäre die
sub-führende $N$-Abhängigkeit der Galerkin-Grundzustände
$\phi^\pm_{\chifn,N}$, kontrolliert durch den Paley-Wiener-Fahrplan aus
\cite{RH_II} adaptiert für Dirichlet $L$-Funktionen (vgl.\
\cite{AnalyticGroundstate}, \S3, und die REG-$\lambda$-konditionierte
Rate aus \cite{RH_II}, \S6). Wir markieren dies als eine offene quantitative
Frage von eigenständigem Interesse.

\begin{remark}[Status der $N=200$-Statistiken]
\label{rem:n200-status}
Die Interpretation des Slope-$0.72$ als $N=200$-spezifische Konstante
hat methodologische Konsequenzen für den gesamten Atlas: Alle bei $N=200$
berichteten Statistiken in Tabellen~\ref{tab:validation},
\ref{tab:regression} und~\ref{tab:archratio} unterliegen demselben
Trunkierungseffekt. Bei Abweichung zwischen $N=600$ und $N=200$ ist der
$N=600$-Wert operationell vorzuziehen. Die $N=200$-Daten werden zur
Transparenz des Konvergenzverhaltens beibehalten, nicht als Endergebnis.
Insbesondere werden die archimedischen Gewichtsverhältnisse $r_\chi$
in Tabelle~\ref{tab:archratio} bei $N=200$ berechnet, weil die
Matrix-Zerlegung $W=W_{\mathrm{arch}}+W_{\mathrm{prime}}$ in dieser
Auflösung am saubersten aufgelöst wird; ihre Werte würden sich bei
$N=600$ moderat verschieben, aber die Prim/Mixed/Arch-Trichotomie
bleibt qualitativ erhalten (Verifikation: vgl.\ Begleitmaterial,
\texttt{ARCH\_TERM\_N600\_ANALYSIS.md}).
\end{remark}

\subsection{Vorzeichenkonvention: algebraische Herleitung des Koeffizienten}
\label{sec:discussion-sign}

Der Koeffizient $-2$ in \eqref{eq:S-def} folgt algebraisch aus der
Struktur der Weil-Quadratform (Gl.~\eqref{eq:weil-qf}): Der primzahlseitige
Beitrag trägt den Faktor $-2\sum_{p,m}w_{p,m}R_\phi(m\log p)$, und der
Sektor-Gap erbt diesen Faktor über die Linearität des Skalarprodukts auf
$H_L^\pm$ und die Sektor-Eigenvektor-Entwicklung. Die vollständige Herleitung,
die alle sechs relevanten Vorzeichenquellen nachverfolgt
(Weil-Formel Primterm-Vorzeichen, Operator-Konvention,
Gap-Orientierung, Paritätsaktion auf Basis, Faltung-vs-Autokorrelation-
Paritätsidentität und der Reell-Charakter-Absorptionsfaktor), wird in
Anhang~\ref{app:sign} und \cite{AnalyticKernelV2}, Theorem~5.1, gegeben.

Die historische Entdeckungssequenz in unseren Entwicklungs-Skripten ---
eine initiale Implementierung mit Koeffizient $+2$, die $R\!=\!-0.70$ auf
der Stichprobe lieferte, korrigiert auf $-2$, die $R\!=\!+0.70$ auf
derselben Stichprobe lieferte --- spiegelt eine frühe Implementierungsfalle
in der Autokorrelations-Paritätskonvention wider (die Vorzeichenrelation
$R_{\phi^-}=-(\phi^-{*}\phi^-)_{\mathrm{math}}$ wurde zunächst falsch
behandelt), kein nachträgliches Fitten des Koeffizienten an die Daten.
Die rigorose Herleitung schließt die Konvention. Die Herleitung liefert
zudem zwei äquivalente Konventionen: die Connes-Operator-Konvention
(Koeffizient $-2$), die im gesamten Atlas verwendet wird, und die
Weil-$Q$-Konvention (Koeffizient $+2$), die sich durch ein Gesamtvorzeichen
am Primterm unterscheiden. Der Atlas hängt nicht von der gewählten Konvention
ab, solange sie konsistent angewendet wird; alle numerischen Vorzeichen sind
konsistent mit der Connes-Operator-Konvention.

\section{Ausblick}
\label{sec:outlook}

Der Atlas liefert eine präzise Diagnose: Eine prädiktive Schließung von
$C^{(2)}_{\chifn}$ über Galerkin-Numerik ist strukturell
unmöglich, und der Weg vorwärts besteht in einer analytischen Kontrolle der
Grundzustands-Differenz, welche Galerkin komplett umgeht. Sechs
konkrete nächste Schritte bieten sich an.

\begin{enumerate}
  \item \textbf{Route~D: $\chifn$-getwisteter CCM-Transfer.}
    Die erste Aufgabe nach Zookeeper ist die Konstruktion des Dirichlet-Analogons
    zum CCM/Zookeeper-Operator: eine $\chifn$-getwistete
    Epstein--Mellin-Transformation, ein Führer-Rang-eins-Term und eine
    $\chifn$-gewichtete arithmetische Perturbation. Die Atlas-Größen
    $G_{\chifn}$, $R_{\chifn}$ und $\eta^\pm_{\chifn}$ sollten dann
    als Diagnostiken für die entsprechende Defektnorm umgelesen werden,
    anstatt als direkte Paley-Wiener-Prädiktoren.
  \item \textbf{Paley-Wiener asymptotische Analyse von
    $\phi^\pm_{\chifn}$.}
    Für den trivialen Charakter ist der asymptotische Grundzustand
    $\phi_\infty$ aus der Paley-Wiener-Maschinerie von \cite{RH_II} (\S5--\S6)
    unter der REG-$\lambda$-Regularitätsannahme explizit bekannt.
    Die Adaption für Dirichlet $L$-Funktionen läuft darauf hinaus, das
    $\chifn$-getwistete Analogon der relevanten Fourier-Multiplikator-Schranke
    und die rationale Funktionskontur zu identifizieren, auf die der
    Residuensatz angewandt wird. Drei der fünf notwendigen Schritte
    lassen sich mechanisch übertragen (vgl.\ \cite{AnalyticGroundstate},
    Thm.~3.1); die verbleibenden zwei erfordern eine charakter-spezifische
    Behandlung von $L'/L(s,\chifn)$ im kritischen Streifen. Dies ist
    die einzelne technische Aufgabe mit der größten Hebelwirkung.
  \item \textbf{Branch-lokaler Hilbert-P\'olya-Rahmen (,,NE-B
    Boundary``).}
    Aufbauend auf \cite{RH_II} (\S4 NE-A, \S5 NE-B) und auf dem
    Selberg-Analogon \cite{Selberg1956} ist das nächste große Projekt im Meta-Programm
    ein strukturelles Theorem, das identifiziert, \emph{welche Zweige des
    Zeta-Funktions-Baums eine Hilbert-P\'olya Interpretation zulassen,
    und welche nicht.} Der vorliegende Atlas liefert die erste
    quantitative Evidenz dafür, dass Dirichlet-Charaktere kleinen Führers
    die zweite Klasse (keine branch-lokale Hilbert-P\'olya-Struktur
    bei endlichem $\lambda$) bevölkern, womit das Randtheorem
    geschärft wird, das den Baum ordnen wird.
  \item \textbf{Führer-Erweiterung $D\leq 100$.}
    Zehn Charaktere stellen einen kleinen Atlas dar. Die Erweiterung auf alle primitiven
    reellen Charaktere mit $D\leq 100$ ($\approx 30$ zusätzliche
    Charaktere, inklusive einiger komplex-konjugierter, die in Paired Spaces
    gruppiert sind) würde statistische Tests der
    Paritäts-/Wurzelzahl-/archimedischen Klassifikation gegen eine
    Verteilung der beobachteten Gap-Vorzeichen ermöglichen. Die Galerkin-Kosten liegen bei
    $\mathcal{O}(D\cdot N^3)$; bei $N=600$ bedeutet dies $\sim 10$ Stunden
    CCX13-Rechenzeit pro zehn Charaktere, also einen Wochenend-Lauf.
  \item \textbf{Paired-Space-Formalismus für komplexe Charaktere.}
    Komplexe Dirichlet-Charaktere treten in konjugierten Paaren auf:
    $\chifn\oplus\bar\chifn$; die Weilsche quadratische Form auf
    $L^2[-L,L]\otimes\mathbb{C}^2$ ist hermitesch, aber nicht
    selbstadjungiert. Der analytische Rahmen aus \S\ref{sec:setup}
    lässt sich verallgemeinern, aber der Sektor-Kern
    (Theorem~\ref{thm:sector-diff}) erhält eine Gaußsummenphase
    (\cite{AnalyticKernel}, Prop.~5.2). Eine saubere Darstellung ist
    die Voraussetzung, um die Wurzelzahl-Achse
    (\S\ref{sec:rootnumber}) gegen nicht-triviale Werte $W\in
    \{+1,-1,\pm i\}$ zu testen.
  \item \textbf{Theoretische Herleitung der Siegel-Walfisz
    Kompression.} Die empirische Dichotomie aus
    Bemerkung~\ref{rem:sw-compression}---zwischen Charakteren mit
    stabilem $\gap_{\mathrm{gal}}^{(N=200)}$ und jenen, die bei $N=600$
    um einen Faktor $\sim 0.35$ komprimiert werden---verlangt nach einer
    theoretischen Klassifikation anhand der charakter-spezifischen
    Siegel-Walfisz Konstante. Eine mögliche Formulierung lautet,
    dass $\chifn$ mit einem echten Bias bei kleinen Primzahlen (Low-Prime Bias, z.\,B.\ $\chifn_{33}$,
    $\chifn_{13}$) eine Primsumme besitzt, welche die
    obere Siegel-Walfisz Schranke im relevanten Fenster ausschöpft und
    daher $N$-stabil ist, wohingegen $\chifn$ mit einer balancierten Primsumme
    (z.\,B.\ $\chifn_{12}$) eine Primsumme unterhalb der oberen
    Schranke aufweist und demzufolge komprimierbar ist.
\end{enumerate}

\appendix

\section{Details zur Berechnung}
\label{app:computation}

\subsection{Galerkin-Basis und Matrixassemblierung}
Wir verwenden die Standard-Fourier-Orthonormalbasis von $H_L^\pm$:
$\{e_0^+=(2L)^{-1/2}\}\cup\{e_n^+(x)=L^{-1/2}\cos(n\pi x/L)\}_{n\geq 1}$
für den geraden Sektor und
$\{e_n^-(x)=L^{-1/2}\sin(n\pi x/L)\}_{n\geq 1}$ für den ungeraden Sektor.
Die Galerkin-trunkierten Operatoren $W^{\pm}_{\chifn,N}$ bei Basisgröße
$N$ werden assembliert durch die Funktion
\texttt{build\_W}$(\chifn,\lambda,\mathrm{sector},N)$ des Skripts
\texttt{arch\_term\_analysis.py}, welche
$W^\pm = W^\pm_{\mathrm{arch}} + W^\pm_{\mathrm{prime}}$ zerlegt, wobei:
\begin{itemize}
\item[(i)] $W^\pm_{\mathrm{arch}}$ aus dem Fourier-Bild
  von $\rho_{\chifn}(t)=\log(q/\pi)+\Re\,\psi(1/4+it/2)$ auf einem Raster
  mit Abstand $\Delta t=10^{-2}$ gewonnen wird, trunkiert bei $|t|\leq T_{\mathrm{arch}}=
  30$ (so gewählt, dass das Spektrum des Gaußschen Grundzustands innerhalb
  der numerischen Genauigkeit enthalten ist).
\item[(ii)] $W^\pm_{\mathrm{prime}}$ durch das Summieren
  von $\chifn(p)^m\log p/p^{m/2}$ über alle $(p,m)$ mit $p\nmid q$ und
  $m\log p\leq 2L$ erhalten wird, wobei die Basisfunktionen-Überlappungen
  (basis-function overlaps) analytisch (für Kosinus-Kosinus) oder durch
  Simpson-Integration (für off-diagonale Sinus-Sinus-Terme) ausgewertet werden.
\end{itemize}
Der kleinste Eigenwert und sein Eigenvektor werden mittels einer
Lanczos-Inverse-Iteration aus \texttt{scipy.sparse.linalg.eigsh}
mit der Toleranz $10^{-10}$ und \texttt{sigma}$=-\log\lambda$ als
Spektralverschiebung (spectral shift) extrahiert.

\subsection{Auswertung der Autokorrelation}
Die Autokorrelation in der Raumdomäne
$(\phi^\pm\ast\phi^\pm)(t)$ wird auf einem uniformen Raster aus $n_{
\mathrm{grid}}=4096$ Punkten auf $[-2L,2L]$ via \texttt{numpy.correlate}
nach Zero-Padding berechnet und dann an den Primzahlabzissen
$t=m\log p$ mittels kubischer \texttt{scipy.interpolate.interp1d} interpoliert.
Die resultierende Realisierung von $S_{\chifn}$, in unseren Skripten
$S_{\mathrm{exact}}$ genannt, unterscheidet sich von der matrix-basierten Realisierung
$\mathrm{prime\_diff}_{\mathrm{matrix}}$ durch Gleitkomma- und
Interpolationsfehler in der Größenordnung von $10^{-3}$ ($1$--$2\%$ von typischen
$|S|$-Werten); dies ist die Ursache für die Vorzeichenwechsel bei $N=600$
für $\chifn_{29}$ und $\chifn_{60}$, die in
\S\ref{sec:validation} dokumentiert sind.

\subsection{Skripte}
Alle Daten in dieser Arbeit wurden erzeugt durch
\texttt{arch\_term\_analysis.py} (für $N=200$ lokal),
\texttt{server\_n600\_full\_analysis.py} (für $N=600$ Server-Lauf,
1.95\,h auf Hetzner CCX13), und
\texttt{chi21\_n\_convergence.py} (für Tabelle~\ref{tab:chi21}).
Die Nachbearbeitung erfolgt durch \texttt{ground\_state\_difference\_analysis.py}.
JSON-Outputs befinden sich in \texttt{CORE/zoo-mapping/\_results/}; Markdown-Zusammenfassungen im selben Verzeichnis.

\section{Status der Vorzeichenkonvention}
\label{app:sign}

Der Koeffizient $-2$ in \eqref{eq:S-def} wird streng abgeleitet aus
der expliziten Weil-Formel in \cite{AnalyticKernelV2},
Theorem~5.1, indem alle sechs relevanten Quellen für Vorzeichen nachverfolgt werden (Minuszeichen der Primzahlseite in der Weil-Formel,
Operatorkonvention, Gap-Orientierung, Wirkung der Parität auf die Basis, Paritätsidentität zwischen
Faltung und Autokorrelation sowie der Absorptionsfaktor für reelle Charaktere).
Die Herleitung liefert \emph{zwei äquivalente Konventionen}:
\begin{enumerate}
  \item die \textbf{Connes-Operator-Konvention} $\tilde A_{\chifn}$,
    in der die Primzahlterme bei der Matrixassemblierung mit positivem Vorzeichen
    addiert werden; dies entspricht dem Skript
    \texttt{build\_W} (d.h.\ $\kappa^{\mathrm{script}}_{\chifn}(u)
    =+\sum w_{p,m}[\delta(u-m\log p)+\delta(u+m\log p)]$) und
    liefert den Koeffizienten $-2$;
  \item die \textbf{Weil-$Q$-Konvention} $Q_{\chifn}$, verwendet in
    \eqref{eq:kappa}, in der $\kappa_{\chifn}$ ein führendes
    Minuszeichen trägt; hier liefert die analoge Herleitung den
    Koeffizienten $+2$ (wie in \cite{AnalyticKernel}, Cor.~4.4).
\end{enumerate}
Die zwei Operatoren erfüllen
$\tilde A_{\chifn}^{\mathrm{prim}}=-Q_{\chifn}^{\mathrm{prim}}$ (die
Primzahl-Teile unterscheiden sich durch ein globales Vorzeichen); beide sind gültige Beschreibungen
desselben spektralen Problems, aber sie führen zu entgegengesetzt orientierten Gaps.
Der Atlas und das Skript verwenden durchgängig konsistent die Connes-Operator-Konvention,
daher ist der Koeffizient $-2$ in
\eqref{eq:S-def} der korrekte.

Eine zweite Subtilität betrifft das Objekt, welches summiert wird. Die numerische
Implementierung von \eqref{eq:S-def} (R1, räumliche Autokorrelation)
verwendet \texttt{np.correlate}, welches die \emph{Autokorrelation}
$R_{\phi^\pm}(u):=\int\phi^\pm(s)\phi^\pm(s-u)\,ds$ berechnet.
Für den \emph{geraden} Grundzustand $\phi^+$ gilt
$R_{\phi^+}=(\phi^+{*}\phi^+)$, doch für den \emph{ungeraden} Grundzustand
$\phi^-$ ergibt die Paritätsidentität
$R_{\phi^-}=-(\phi^-{*}\phi^-)_{\mathrm{math}}$ (Beweis: Substitution $v=-s$ im
Integral $\int\phi^-(s)\phi^-(s-u)\,ds$, unter Nutzung von
$\phi^-(-v)=-\phi^-(v)$). Folglich ist der
numerische Ausdruck
$\Deltachi^{\mathrm{num}}:=R_{\phi^+}-R_{\phi^-}$, der vom Skript berechnet wird, gleich
$(\phi^+{*}\phi^+)_{\mathrm{math}}+(\phi^-{*}\phi^-)_{\mathrm{math}}$ in Bezug auf mathematische
Faltungen --- also eine \emph{Summe}, die genau der rechten Seite von~\eqref{eq:Delta-def}
entspricht. Damit stimmt $\Deltachi$ wie in~\eqref{eq:Delta-def} definiert identisch mit der
skript-berechneten Größe $\Deltachi^{\mathrm{num}}$ überein. Der Koeffizient~$-2$ in
\eqref{eq:S-def} ist dann der eindeutige Wert, der die Gap-Identität
(Satz~\ref{thm:closure}) erfüllt. Siehe
\cite{AnalyticKernelV2}, \S4 und Theorem~5.1, für die vollständige
Herleitung. Alle gemeldeten Vorzeichen und sämtliche Tabellen
(Tabellen~\ref{tab:validation}, \ref{tab:regression},
\ref{tab:archratio}, \ref{tab:chi21}) sind in sich konsistent
und wurden gegen eine direkte Galerkin-Diagonalisierung verifiziert.

% --- Bibliografie ---
\begin{thebibliography}{99}

% ----- FST Master-Referenzen -----

\bibitem{MathMaster}
L.~Geiger,
\textit{The Zeta Zoo --- The Mathematical Side of Functional
Stability Theory: Branch-Local Hilbert--P\'olya via the Trinity of
UCU, SGE, and Weil-Form Transversality},
Preprint, Zenodo (2026), DOI:
\href{https://doi.org/10.5281/zenodo.19673226}{10.5281/zenodo.19673226}.
Dreikomponenten-Zerlegung von $\mathrm{RH}$ in \S2.4; UCU in \S7;
CRM-V Klassifikations-Bemerkung in \S6.5.

\bibitem{RFEPFramework}
L.~Geiger,
\textit{FST Spectrum Duality: The Renormalized Free Energy Principle
as Physical Instantiation of Functional Stability Theory},
Preprint, Zenodo (2026), DOI:
\href{https://doi.org/10.5281/zenodo.19036190}{10.5281/zenodo.19036190}.

\bibitem{Zookeeper2026}
L.~Geiger,
\textit{The Spectral Zookeeper: The Riemann Hypothesis via Spectral
Microcluster Closure in the Connes--Consani--Moscovici Framework},
Preprint (2026), Zenodo DOI
\href{https://doi.org/10.5281/zenodo.19673126}{10.5281/zenodo.19673126}.

% ----- Interne Projektreferenzen (Preprints / Ergänzungsnotizen) -----

\bibitem{RH_II}
L.~Geiger,
\textit{A Conditional Reduction of the Riemann Hypothesis to Even
Dominance of the Weil Quadratic Form},
Preprint (2026), Zenodo-Concept-DOI
\href{https://doi.org/10.5281/zenodo.19764771}{10.5281/zenodo.19764771}.
Zuletzt geprüfte Version: v1.8, Record-DOI
\href{https://doi.org/10.5281/zenodo.20396370}{10.5281/zenodo.20396370}.

\bibitem{MetaGalerkinBridge}
L.~Geiger,
\textit{A Meta Galerkin Bridge: Diagnostic Decomposition of RH-Type
Programmes}
(Entwurf 2026), archivierter Vorgänger des \texttt{CORE/zetazoo}
Boundary-Papers. Inhalt absorbiert in \cite{MathMaster} \S2.4.

\bibitem{AnalyticKernel}
L.~Geiger,
\textit{Analytic Kernel: Sector-Kernel Decomposition of the
Dirichlet Weil Quadratic Form},
Ergänzungsnotiz (2026-04-17, revidiert 2026-04-16),
\path{CORE/zoo-mapping/ANALYTIC_PIPELINE.md}.

\bibitem{AnalyticKernelV2}
L.~Geiger,
\textit{Analytic Kernel V2 --- Rigorous Weil-Formula Derivation
of the Sign Coefficient},
Ergänzungsnotiz (2026-04-17),
\path{CORE/zoo-mapping/ANALYTIC_PIPELINE.md}.
Löst den Status der Vorzeichenkonvention durch stringente Ableitung des
Koeffizienten aus der expliziten Weil-Formel unter Berücksichtigung aller sechs Vorzeichenquellen.

\bibitem{AnalyticGroundstate}
L.~Geiger,
\textit{Analytic Ground State: Paley-Wiener Asymptotics for the
Sector Ground States of Dirichlet Weil Quadratic Forms},
Ergänzungsnotiz (2026-04-17),
\path{CORE/zoo-mapping/ANALYTIC_PIPELINE.md}.

\bibitem{Woche4Validation}
L.~Geiger,
\textit{Week-4 Validation: 16-Window Falsification of the
Leading-Order Gap Formula},
Ergänzungsnotiz (2026-04-17),
\path{CORE/zoo-mapping/_results/WOCHE4_VALIDATION_RESULT.md}.

\bibitem{AsymptoticScan}
L.~Geiger,
\textit{Asymptotic Scan of the $\chi_5$ and $\chi_{12}$ Sector Gaps
across Three Orders of Magnitude in $\lambda$},
Ergänzungsnotiz (2026-04-15),
\path{CORE/zoo-mapping/_results/ASYMPTOTIC_SCAN_2026-04-15.md}.

% ----- Externe Referenzen (überprüft 2026-05-15) -----

\bibitem{IwaniecKowalski}
H.~Iwaniec und E.~Kowalski,
\textit{Analytic Number Theory},
American Mathematical Society Colloquium Publications \textbf{53},
AMS, Providence, RI (2004), ISBN 0-8218-3633-1, 615~S.

\bibitem{Weil1952}
A.~Weil,
\textit{Sur les ``formules explicites'' de la th\'eorie des nombres
premiers},
Meddelanden fr\aa n Lunds Universitets Matematiska Seminarium
(Vol.\ ded.\ \`a Marcel Riesz), 252--265 (1952).

\bibitem{Selberg1956}
A.~Selberg,
\textit{Harmonic Analysis and Discontinuous Groups in Weakly
Symmetric Riemannian Spaces with Applications to Dirichlet Series},
J. Indian Math. Soc. \textbf{20}, 47--87 (1956).

\bibitem{Connes}
A.~Connes,
\textit{Trace formula in noncommutative geometry and the zeros of
the Riemann zeta function},
Selecta Math. (New Series) \textbf{5}, 29--106 (1999),
DOI \href{https://doi.org/10.1007/s000290050042}{10.1007/s000290050042}.

\bibitem{PaleyWiener}
R.~E.~A.~C.~Paley und N.~Wiener,
\textit{Fourier Transforms in the Complex Domain},
American Mathematical Society Colloquium Publications \textbf{19},
AMS, Providence, RI (1934).

\bibitem{LMFDB}
The LMFDB Collaboration,
\textit{The L-Functions and Modular Forms Database},
\url{https://www.lmfdb.org/}. L-Funktionsseiten abgerufen am 16.05.2026.

\end{thebibliography}

\end{document}
